% COMPLEX ANALYSIS -- THE RISING SEA FORMAT
% MATH 503, Fall 2025, Lectures 1--21. Proofs intentionally omitted.
% Compile twice with XeLaTeX or LuaLaTeX.
\RequirePackage{silence}
\WarningFilter{xcolor}{Package option}
\documentclass[10pt,letterpaper,twoside,openright]{amsbook}
% TYPOGRAPHY AND MATHEMATICS
\usepackage[no-math]{fontspec}
\IfFontExistsTF{TeX Gyre Pagella}
 {\newfontfamily\SeaPhraseFont{TeX Gyre Pagella}}
 {\newfontfamily\SeaPhraseFont{Latin Modern Roman}}
\usepackage[T1]{fontenc}
\usepackage{palatino}
\renewcommand{\updefault}{n}
\usepackage{amsmath,amsthm,amssymb,mathtools}
\usepackage{euler}
\usepackage{mathrsfs}
\usepackage{microtype}
% SMALL CAPS AND LEAD-IN PHRASES
\providecommand{\smallcaps}[1]{}
\providecommand{\allcaps}[1]{}
\renewcommand{\smallcaps}[1]{{\addfontfeatures{Letters=SmallCaps,LetterSpace=4}#1}}
\renewcommand{\allcaps}[1]{{\addfontfeatures{LetterSpace=8}\MakeUppercase{#1}}}
\newcommand{\leadphrase}[1]{{\addfontfeatures{Letters=SmallCaps,LetterSpace=4}#1}}
\AddToHook{cmd/smallcaps/before}{\begingroup\SeaPhraseFont}
\AddToHook{cmd/smallcaps/after}{\endgroup}
\AddToHook{cmd/allcaps/before}{\begingroup\SeaPhraseFont}
\AddToHook{cmd/allcaps/after}{\endgroup}
\AddToHook{cmd/leadphrase/before}{\begingroup\SeaPhraseFont}
\AddToHook{cmd/leadphrase/after}{\endgroup}
% COLORS
\usepackage{xcolor}
\definecolor{PageBlack}{HTML}{FFFFFF}
\definecolor{SoftWhite}{HTML}{464646}
\definecolor{MathYellow}{HTML}{003496}
\definecolor{SoftGray}{HTML}{000000}
\definecolor{SeaLink}{HTML}{000080}
\definecolor{SeaCite}{HTML}{008000}
% PAGE GEOMETRY
\setlength{\textwidth}{30pc}
\setlength{\textheight}{50.5pc}
\setlength{\headheight}{8pt}
\setlength{\headsep}{14pt}
\setlength{\footskip}{18pt}
\setlength{\topskip}{10pt}
\calclayout
\setlength{\normalparindent}{18pt}
\setlength{\parindent}{18pt}
\setlength{\parskip}{0pt}
\setlength{\emergencystretch}{2em}
\flushbottom
% DIAGRAMS, LISTS, AND TABLES
\usepackage{tikz}
\usepackage{tikz-cd}
\usetikzlibrary{arrows.meta,calc,positioning,matrix,decorations.pathmorphing}
\tikzset{every picture/.style={color=SoftWhite},
 every node/.style={text=SoftWhite},every path/.style={draw=SoftWhite}}
\usepackage{booktabs,array,graphicx,enumitem,longtable}
\setlist[itemize]{leftmargin=1.8em,itemsep=0pt,topsep=\smallskipamount,parsep=0pt}
\setlist[enumerate]{label=(\alph*),leftmargin=1.8em,itemsep=0pt,topsep=\smallskipamount,parsep=0pt}
% NUMBERING, HEADINGS, AND TABLE OF CONTENTS
\setcounter{secnumdepth}{2}
\setcounter{tocdepth}{1}
\renewcommand{\thepart}{\Roman{part}}
\renewcommand{\thesection}{\thechapter.\arabic{section}}
\renewcommand{\thesubsection}{\thesection.\arabic{subsection}}
\numberwithin{equation}{subsection}
\makeatletter
\renewcommand{\section}{\@startsection{section}{1}
 {\z@}{2\baselineskip plus .5\baselineskip}
 {\baselineskip}{\normalfont\Large\bfseries\centering}}
\renewcommand{\subsection}{\@startsection{subsection}{2}
 {\z@}{\medskipamount}{-.5em}{\normalfont\bfseries}}
\renewcommand{\@seccntformat}[1]{%
 \csname the#1\endcsname
 \ifnum\csname c@secnumdepth\endcsname>\z@\relax
 \expandafter\ifx\csname #1\endcsname\section \enspace
 \else .\enspace \fi\fi}
\renewcommand{\frontmatter}{\cleardoublepage\pagenumbering{arabic}}
\renewcommand{\mainmatter}{\cleardoublepage}
\renewcommand{\tocsection}[3]{%
 \indentlabel{\@ifnotempty{#2}{\ignorespaces#1 #2.\quad\hspace{3.11pt}}}#3}
\makeatother
% THEOREM ENVIRONMENTS
\newtheoremstyle{sea-plain}
 {\medskipamount}{\medskipamount}{\itshape}{0pt}{\bfseries}{}{.5em}
 {\thmnumber{#2. }\thmname{#1}\thmnote{ (#3)}.\ ---}
\newtheoremstyle{sea-definition}
 {\medskipamount}{\medskipamount}{\normalfont}{0pt}{\normalfont}{}{.5em}
 {\thmnumber{\textbf{#2.} }\thmname{\textit{#1\thmnote{: #3}.}}}
\newtheoremstyle{sea-exercise}
 {\medskipamount}{\medskipamount}{\normalfont}{0pt}{\normalfont}{}{.5em}
 {\thmnumber{\textbf{#2.} }\thmname{\textsc{#1}}\thmnote{ (#3)}.}
\theoremstyle{sea-plain}
\newtheorem{theorem}[subsection]{Theorem}
\newtheorem{proposition}[subsection]{Proposition}
\newtheorem{lemma}[subsection]{Lemma}
\newtheorem{corollary}[subsection]{Corollary}
\newtheorem{claim}[subsection]{Claim}
\theoremstyle{sea-definition}
\newtheorem{definition}[subsection]{Definition}
\newtheorem{example}[subsection]{Example}
\newtheorem{construction}[subsection]{Construction}
\newtheorem{remark}[subsection]{Remark}
\newtheorem{note}[subsection]{Note}
\newtheorem{warning}[subsection]{Warning}
\newtheorem{conjecture}[subsection]{Conjecture}
\theoremstyle{sea-exercise}
\newtheorem{exercise}{Exercise}[section]
\newtheorem{problem}[subsection]{Problem}
\usepackage{alphalph}
\renewcommand{\theexercise}{\thesection.\AlphAlph{\value{exercise}}}
\makeatletter
\renewenvironment{proof}[1][\proofname]{%
 \par\pushQED{\qed}\normalfont
 \topsep6\p@\@plus6\p@\relax\trivlist\itemindent\z@
 \item[\hskip\labelsep\itshape#1\@addpunct{.}]\ignorespaces
}{\popQED\endtrivlist\@endpefalse}
\makeatother
\newenvironment{solution}{\begin{proof}[Solution]}{\end{proof}}
% OPERATORS AND NOTATION FOR COMPLEX ANALYSIS
\DeclareMathOperator{\Aut}{Aut}
\DeclareMathOperator{\im}{im}
\DeclareMathOperator{\Res}{Res}
\DeclareMathOperator{\Ind}{Ind}
\DeclareMathOperator{\Log}{Log}
\DeclareMathOperator{\Arg}{Arg}
\DeclareMathOperator{\Real}{Re}
\DeclareMathOperator{\Imag}{Im}
\DeclareMathOperator{\diam}{diam}
\DeclareMathOperator{\dist}{dist}
\DeclareMathOperator{\supp}{supp}
\DeclareMathOperator{\sgn}{sgn}
\newcommand{\NN}{\mathbb{N}}
\newcommand{\ZZ}{\mathbb{Z}}
\newcommand{\QQ}{\mathbb{Q}}
\newcommand{\RR}{\mathbb{R}}
\newcommand{\CC}{\mathbb{C}}
\newcommand{\DD}{\mathbb{D}}
\newcommand{\HH}{\mathbb{H}}
\newcommand{\Chat}{\mathbb{C}_{\infty}}
\newcommand{\Hol}{H}
\newcommand{\Sch}{\mathcal{S}}
\newcommand{\Fcal}{\mathcal{F}}
\newcommand{\into}{\hookrightarrow}
\newcommand{\onto}{\twoheadrightarrow}
\newcommand{\isom}{\xrightarrow{\sim}}
\newcommand{\dd}{\,d}
\newcommand{\bD}[2]{\overline{D(#1,#2)}}
\newcommand{\pD}[2]{D'(#1,#2)}
\newcommand{\norm}[1]{\left\lVert#1\right\rVert}
\definecolor{SplitLeft}{HTML}{B13A34}
\definecolor{SplitRight}{HTML}{1F5FA8}
\definecolor{OldInterval}{HTML}{803E9B}
\newcommand{\leftpiece}[1]{\textcolor{SplitLeft}{#1}}
\newcommand{\rightpiece}[1]{\textcolor{SplitRight}{#1}}
\newcommand{\oldpiece}[1]{\textcolor{OldInterval}{#1}}
\newcommand{\lectureinfo}[2]{%
 \par\noindent{\small\textit{MATH 503, Fall 2025. #1. Source PDF pages #2.}}\par\medskip}
\newcommand{\sourcefigure}[3]{%
 \begin{figure}[htbp]\centering
 \includegraphics[width=#2\linewidth]{figures/#1.pdf}
 \caption{#3}\end{figure}}
% TITLE, RUNNING HEADS, AND HYPERLINKS
\title{Complex Analysis}
\newcommand{\booksubtitle}{Lecture Notes without Proofs}
\author{Notes from MATH 503, Fall 2025}
\date{}
\newcommand{\evenrunninghead}{Complex Analysis}
\newcommand{\oddrunninghead}{MATH 503 Lecture Notes}
\makeatletter
\renewcommand{\maketitle}{%
 \begin{titlepage}\setcounter{page}{1}\thispagestyle{empty}\centering
 \vspace*{-13.8pt}
 {\fontsize{20.74}{25}\selectfont\bfseries\MakeUppercase{\@title}\par}
 \vspace{133.7pt}
 {\fontsize{17.28}{22}\selectfont\bfseries\booksubtitle\par}
 \vfill {\small\authors\par}
 \ifx\@date\@empty\else\medskip{\small\@date\par}\fi
 \vspace*{72pt}
 \end{titlepage}}
\def\ps@sea{%
 \ps@empty
 \def\@evenhead{\normalfont\scriptsize\rlap{\thepage}\hfil\evenrunninghead\hfil}%
 \def\@oddhead{\normalfont\scriptsize\hfil\oddrunninghead\hfil\llap{\thepage}}}
\makeatother
\usepackage{hyperref}
\hypersetup{
 colorlinks=true,linkcolor=SeaLink,citecolor=SeaCite,urlcolor=SeaLink,
 linktoc=page,pdfborder={0 0 0},bookmarksdepth=section,
 pdftitle={Complex Analysis: MATH 503 Lecture Notes},
 pdfauthor={MATH 503, Fall 2025}}
\pdfstringdefDisableCommands{%
 \def\enspace{ }\def\noindent{}\def\smallcaps#1{#1}\def\allcaps#1{#1}\def\leadphrase#1{#1}
 \def\CC{C}\def\RR{R}\def\DD{D}\def\HH{H}\def\Chat{C infinity}}

\newcommand{\src}[1]{\par\smallskip\noindent
 {\scriptsize\sffamily\color{SoftWhite}Source slide #1.}\par\nopagebreak\smallskip}
\newcommand{\lecturedetails}[2]{\par\noindent{\small
 MATH 503, Fall 2025. #1. Merged PDF pages #2.}\par\medskip}
\newcommand{\editorial}[1]{\par\smallskip\noindent
 {\small\textit{Editorial clarification.} #1}\par\smallskip}
\begin{document}
\pagecolor{PageBlack}\color{SoftWhite}
\everymath{\color{MathYellow}}
\AtBeginEnvironment{equation}{\color{MathYellow}}
\AtBeginEnvironment{equation*}{\color{MathYellow}}
\AtBeginEnvironment{align}{\color{MathYellow}}
\AtBeginEnvironment{align*}{\color{MathYellow}}
\frontmatter\pagestyle{plain}\maketitle
\chapter*{Reading this edition}
This edition follows the source order of the twenty-one lectures in
\emph{Complex Analysis Merged Notes(2).pdf}, MATH 503, Fall 2025.
Source-slide numbers refer to pages of that merged PDF, not to the page
numbers of this book. Definitions, statements, remarks, explanatory prose,
and examples are kept; proofs are omitted. Explicit \emph{Prove it},
\emph{Why?}, and similar reader prompts are recorded as exercises.

The typesetting and numbering belong to this edition. Original
illustrations are reproduced as vector figures. Recalled results are retained
where they occur again. Editorial clarifications are identified separately.
The accompanying coverage checklist records how every source page was handled.
\tableofcontents
\mainmatter\pagestyle{sea}
\chapter{Complex Numbers}
\lecturedetails{September 4, 2025}{1--27}
\section{The complex field}
\src{2}
\begin{definition}[Complex numbers]
A complex number is an ordered pair $(a,b)\in\RR\times\RR$.
\end{definition}
\begin{definition}[Addition and multiplication of complex numbers]
For two complex numbers $x=(a,b)$ and $y=(c,d)$, addition and multiplication
are defined by
\[
 x+y=(a+c,b+d),\qquad x\cdot y=(ac-bd,ad+bc).
\]
\end{definition}
\src{3}
\begin{theorem}[Complex field]
These operations turn the set of all complex numbers into a field, with
$(0,0)$ and $(1,0)$ playing respectively the roles of $0$ and $1$.
This field is denoted by $\CC$.
\end{theorem}
\src{8}
\begin{remark}
For any $a,b\in\RR$,
\[
 (a,0)+(b,0)=(a+b,0),\qquad (a,0)(b,0)=(ab,0).
\]
The complex numbers in $\{(a,0):a\in\RR\}$ have the same arithmetic
properties as the corresponding real numbers. We can therefore identify
$(a,0)$ with $a$. This identification gives the real field $\RR$ as a
subfield of $\CC$.

We have defined $\CC$ without reference to the mysterious square root of
$-1$. We now show that $(a,b)$ is equivalent to the customary $a+bi$.
\end{remark}
\begin{definition}
The imaginary number is $i=(0,1)$.
\end{definition}
\src{9}
\begin{theorem}
One has $i^2=-1$.
\end{theorem}
\begin{theorem}
One has $\CC=\{a+ib:a,b\in\RR\}$.
\end{theorem}
\src{10}
\begin{definition}[Conjugate, real and imaginary parts]
For $z=a+ib$ with $a,b\in\RR$, the number $\bar z=a-ib$ is its conjugate.
The numbers $a,b$ are respectively its real and imaginary parts:
\[
 a=\Re(z)=\Real(z),\qquad b=\Im(z)=\Imag(z).
\]
\end{definition}
\begin{theorem}
For $z,w\in\CC$:
\begin{enumerate}[label=(\roman*)]
\item $\overline{z+w}=\bar z+\bar w$.
\item $\overline{zw}=\bar z\,\bar w$.
\item $z+\bar z=2\Real z$ and $z-\bar z=2i\Imag z$.
\item $z\bar z$ is a positive real number except when $z=0$.
\end{enumerate}
\end{theorem}
\src{12}
\begin{definition}[Absolute value]
For $z\in\CC$, its absolute value is $|z|=\sqrt{z\bar z}$.
\end{definition}
\begin{remark}
This absolute value exists and is unique. It coincides with the real
absolute value: if $x\in\RR$, then $\bar x=x$, and
\[
 |x|=\sqrt{x\bar x}=\sqrt{x^2}
 =\begin{cases}x,&x\geq0,\\-x,&x<0.\end{cases}
\]
\end{remark}
\src{13}
\begin{theorem}[Properties of absolute value]
For $z,w\in\CC$:
\begin{enumerate}[label=(\roman*)]
\item $|z|>0$ if and only if $z\ne0$, and $|0|=0$.
\item $|\bar z|=|z|$.
\item $|zw|=|z||w|$.
\item $|\Real z|\leq|z|$ and $|\Imag z|\leq|z|$.
\item $|z+w|\leq|z|+|w|$.
\end{enumerate}
\end{theorem}
\src{15}
\begin{definition}[Argument]
An argument $\arg z$ of $z=a+ib$ is the angle which the segment from
$(0,0)$ to $(a,b)$ makes with the positive real axis. The argument is not
unique but is determined up to a multiple of $2\pi$.
\end{definition}
If $r=|z|$ and $\theta=\arg z$, then
\[
 z=r(\cos\theta+i\sin\theta).
\]
Trigonometric identities give
\[
 \arg(z_1z_2)=\arg z_1+\arg z_2,\qquad
 \arg(z_1/z_2)=\arg z_1-\arg z_2,\qquad
 \arg\bar z=-\arg z.
\]
An argument is \emph{principal} if it lies in $(-\pi,\pi]$.
\editorial{Arguments are defined only for nonzero numbers. The displayed
argument equalities are understood modulo $2\pi$.}
\section{Topology of the complex plane}
\src{16}
\begin{definition}[Convergence]
A sequence $(z_n)_{n\in\NN}\subseteq\CC$ converges to $z$, written
$\lim_{n\to\infty}z_n=z$, if and only if $\lim_{n\to\infty}|z_n-z|=0$.
Equivalently, for every $\varepsilon>0$ there is $N_\varepsilon\in\NN$ such
that $n\geq N_\varepsilon$ implies $|z_n-z|<\varepsilon$.
\end{definition}
One has $z_n\to z$ if and only if
$\Real z_n\to\Real z$ and $\Imag z_n\to\Imag z$.
We write $z_n\to\infty$ when $|z_n|\to\infty$.
\src{17}
\begin{definition}[Cauchy sequence]
A sequence $(z_n)$ is Cauchy if
$\lim_{m,n\to\infty}|z_n-z_m|=0$. Equivalently, for every $\varepsilon>0$
there is $N_\varepsilon$ such that $m,n\geq N_\varepsilon$ implies
$|z_n-z_m|<\varepsilon$.
\end{definition}
\begin{theorem}
The complex plane with metric $d(z_1,z_2)=|z_1-z_2|$ is a complete metric space.
\end{theorem}
\src{18}
\begin{definition}[Discs, punctured discs, and circles]
For $a\in\CC$ and $r>0$,
\[
\begin{aligned}
D(a,r)&=\{z:|z-a|<r\},&
\overline{D(a,r)}&=\{z:|z-a|\leq r\},\\
D'(a,r)&=\{z:0<|z-a|<r\},&
C(a,r)&=\{z:|z-a|=r\}.
\end{aligned}
\]
The unit disc is $\DD=D(0,1)$. The punctured disc is open, and
$C(a,r)=\overline{D(a,r)}\setminus D(a,r)$.
\end{definition}
\src{19}
A set $\Omega\subseteq\CC$ is \emph{open} if each $a\in\Omega$ has a disc
$D(a,r)\subseteq\Omega$. It is \emph{closed} if its complement
$\Omega^c=\CC\setminus\Omega$ is open.

A point $z$ is a \emph{limit point} of $\Omega$ if there is a sequence
$(z_n)\subseteq\Omega$ with $z_n\ne z$ and $z_n\to z$.
A set is closed if and only if it contains all its limit points. Its closure
$\overline\Omega$ is the union of the set and its limit points.
The boundary $\partial\Omega$ is its closure minus its interior.
For instance, $C(a,r)=\partial D(a,r)$.
\src{20}
For $a,b\in\CC$, $[a,b]$ denotes the closed segment with endpoints $a,b$.
For arbitrary $t_1<t_2$,
\[
 [a,b]=\left\{a+\frac{t-t_1}{t_2-t_1}(b-a):t_1\leq t\leq t_2\right\}.
\]
For $a_1,\ldots,a_{n+1}\in\CC$, a \emph{polygon} joining $a_1$ to
$a_{n+1}$ is $\bigcup_{j=1}^n[a_j,a_{j+1}]$, abbreviated
$[a_1,\ldots,a_{n+1}]$.
\src{21}
\begin{definition}[Connectedness]
Let $(X,d)$ be a metric space and $E\subseteq X$. The set $E$ is connected
if it cannot be written as a disjoint union of two nonempty relatively open
subsets of $E$.
\end{definition}
Thus $E=A\cup B$, $A\cap B=\varnothing$, with $A,B$ open in $E$, implies
$A=\varnothing$ or $B=\varnothing$. Otherwise this is called a
\emph{separation} of $E$ into open sets. The union of two disjoint open
discs $A,B$ is not connected, since
\[
 E=A\cup B=(A\cap E)\cup(B\cap E),
\]
where the two sets are nonempty, disjoint, and relatively open.
An open connected set is called a \emph{region}.
\src{22}
\begin{definition}[Component]
A maximal connected subset of $E$ is a component of $E$.
\end{definition}
For $a\in E$, let $C(a)$ be the union of all connected subsets of $E$
containing $a$. Since $\{a\}$ is connected, $a\in C(a)$, and
$E=\bigcup_{a\in E}C(a)$.
\begin{lemma}
\begin{enumerate}[label=(\roman*)]
\item $C(a)$ is connected.
\item Components of $E$ are disjoint or identical.
\item Components of an open set are open.
\end{enumerate}
\end{lemma}
\begin{theorem}
An open set in the complex plane is a disjoint union of regions.
\end{theorem}
\editorial{The slide states the last two openness assertions for arbitrary
metric spaces. They require local connectedness; here they are applied to
$\CC$, where they hold.}
\src{25}
\begin{theorem}[Polygonal connectedness]
A nonempty open subset $E$ of $\CC$ is connected if and only if any two
points of $E$ can be joined by a polygonal path lying in $E$.
\end{theorem}

\chapter{Compactification, Holomorphic Functions, and Cauchy--Riemann Equations}
\lecturedetails{September 8, 2025}{28--59}
\section{The Riemann sphere}
\src{29}
Let $(X,d)$ be a metric space and $\Omega\subseteq X$. An open covering of
$\Omega$ is a family $(U_\alpha)_{\alpha\in A}$ of open sets such that
$\Omega\subseteq\bigcup_{\alpha\in A}U_\alpha$. The index set need not be countable.
\begin{definition}[Compactness]
A set in a metric space is compact if and only if every open covering has
a finite subcovering.
\end{definition}
\begin{theorem}
A set $\Omega$ in a metric space is compact if and only if every sequence
in $\Omega$ has a subsequence converging to a point of $\Omega$.
\end{theorem}
\begin{theorem}
A subset of $\CC$ is compact if and only if it is closed and bounded.
\end{theorem}
\src{30}
Adjoining a point $\infty$, called the point at infinity, gives the
\emph{extended complex plane} $\Chat=\CC\cup\{\infty\}$.
The Riemann sphere is
\[
 S=\{(x_1,x_2,x_3)\in\RR^3:x_1^2+x_2^2+x_3^2=1\}.
\]
The \emph{stereographic projection} $f:S\to\Chat$ is
\[
 f(x_1,x_2,x_3)=\frac{x_1+ix_2}{1-x_3}\quad
 ((x_1,x_2,x_3)\ne(0,0,1)),\qquad f(0,0,1)=\infty.
\]
The point $N=(0,0,1)$ is the north pole. The map is one-to-one and onto.
\src{33--35}
Identify $z=x+iy$ with $(x,y,0)\in\RR^3$. The line joining $z$ to $N$ is
\[
 tN+(1-t)z=((1-t)x,(1-t)y,t),\qquad-\infty<t<\infty.
\]
Its intersections with $S$ satisfy
\[
 (1-t)^2x^2+(1-t)^2y^2+t^2=1.
\]
Away from the north pole, this becomes
$(1-t)|z|^2=1+t$, hence $t=(|z|^2-1)/(|z|^2+1)$.
\sourcefigure{p034}{.9}{The stereographic projection diagram from source slide 34.}
The other intersection is
\[
\begin{aligned}
 Z&=\left(\left[1-\frac{|z|^2-1}{|z|^2+1}\right]x,
 \left[1-\frac{|z|^2-1}{|z|^2+1}\right]y,
 \frac{|z|^2-1}{|z|^2+1}\right)\\
 &=\left(\frac{2x}{|z|^2+1},\frac{2y}{|z|^2+1},
 \frac{|z|^2-1}{|z|^2+1}\right)\\
 &=\left(\frac{z+\bar z}{|z|^2+1},
 \frac{z-\bar z}{i(|z|^2+1)},\frac{|z|^2-1}{|z|^2+1}\right).
\end{aligned}
\]
Thus the line through $(x,y,0)$ and $(0,0,1)$ meets $S\setminus\{N\}$ in
exactly one point $Z=(x_1,x_2,x_3)$, and $f(Z)=z$.
These are the \emph{spherical coordinates} of $z$; those of $\infty$ are
$(0,0,1)$.
\src{36--38}
For $z,z'\in\Chat$, let $x=(x_1,x_2,x_3)$ and $x'=(x'_1,x'_2,x'_3)$
be their spherical coordinates, as given above for each finite point. Define
\[
 d(z,z')=\sqrt{(x_1-x'_1)^2+(x_2-x'_2)^2+(x_3-x'_3)^2}.
\]
Consequently
\[
 d(z,z')^2=2-2(x_1x'_1+x_2x'_2+x_3x'_3).
\]
Writing $D=(|z|^2+1)(|z'|^2+1)$, the scalar product calculation is
\[
\begin{aligned}
 x_1x'_1+x_2x'_2+x_3x'_3
 &=\frac{(z+\bar z)(z'+\bar z')-(z-\bar z)(z'-\bar z')
       +( |z|^2-1)(|z'|^2-1)}{D}\\
 &=\frac{(|z|^2-1)(|z'|^2-1)+2z\bar z'+2\bar z z'}{D}\\
 &=\frac{D-2(|z|^2+1)-2(|z'|^2+1)+4+2z\bar z'+2\bar z z'}{D}.
\end{aligned}
\]
The numerator subtracted from $D$ is
\[
\begin{aligned}
 &2(|z|^2+1)+2(|z'|^2+1)-4-2z\bar z'-2\bar z z'\\
 &\qquad=2(|z|^2+|z'|^2-z\bar z'-\bar z z')=2|z-z'|^2.
\end{aligned}
\]
Therefore
\[
 2(x_1x'_1+x_2x'_2+x_3x'_3)=2-\frac{4|z-z'|^2}{D},
\quad d(z,z')^2=\frac{4|z-z'|^2}{D},
\]
\[
 d(z,z')=\frac{2|z-z'|}{\sqrt{(|z|^2+1)(|z'|^2+1)}}.
\]
If $z'=\infty$, then $x'_1=x'_2=0$, $x'_3=1$, and
\[
 d(z,\infty)^2=2-2x_3
 =2-2\frac{|z|^2-1}{|z|^2+1}=\frac4{|z|^2+1}.
\]
\src{39}
Hence $d(z,\infty)=2/\sqrt{|z|^2+1}$, and $\Chat$ is a metric space.
For $E\subseteq\Chat$, the induced topology is characterized as follows:
\begin{itemize}
\item If $\infty\notin E$, then $E$ is open in $\Chat$ if and only if it is
open in $\CC$.
\item If $\infty\in E$, then $E$ is open if and only if its complement in
$\Chat$ is a closed bounded (equivalently compact) subset of $\CC$.
\end{itemize}
Checking continuity of $f$ and $f^{-1}$ gives:
\begin{theorem}
The extended complex plane is homeomorphic to $S$. In particular, it is compact.
\end{theorem}
\section{Complex functions}
\src{40}
A function $f$ on $E\subseteq\CC$ assigns to each $z\in E$ a unique complex
number $w=f(z)$, its value at $z$. The set $E$ is the domain of definition.
If $z=x+iy$ and $w=u+iv=f(x+iy)$, then $u,v$ depend on the real variables
$x,y$, and
\[
 f(z)=u(x,y)+iv(x,y),\qquad u=\Real f,\quad v=\Imag f.
\]
\src{41}
\begin{example}
For $f(z)=z^2$, $f(x+iy)=x^2-y^2+2ixy$, so
$\Real f=x^2-y^2$ and $\Imag f=2xy$. Its domain is all of $\CC$.
\end{example}
\begin{example}[Branches of roots]
Assigning all $n$th roots to each complex number does not give a function
in the preceding single-valued sense: each $z\ne0$ has $n$ roots.
If $z=|z|(\cos\theta+i\sin\theta)$ with $\theta\in(-\pi,\pi]$, they are
\[
 z_k=\sqrt[n]{|z|}
 \left(\cos\frac{\theta+2k\pi}{n}
 +i\sin\frac{\theta+2k\pi}{n}\right),\qquad k=0,\ldots,n-1.
\]
These determine $n$ branches. Choosing a particular $k$ makes a
single-valued function. The $k=0$ branch is the principal branch.
\end{example}
\src{42}
\begin{definition}[Continuity]
For $f:E\to\CC$ and $z_0\in E$, continuity at $z_0$ means that, for every
$\varepsilon>0$, there is $\delta>0$ such that
$z\in E$ and $|z-z_0|<\delta$ imply $|f(z)-f(z_0)|<\varepsilon$.
Equivalently, whenever $z_n\in E$ and $z_n\to z_0$, one has
$f(z_n)\to f(z_0)$.
\end{definition}
Continuity on $E$ means continuity at every point. Sums and products of
continuous functions are continuous. A complex function is continuous if
and only if both its real and imaginary parts are continuous.
\src{43}
Convergence in $\CC$ and in $\RR^2$ is the same, so continuity as a complex
function is equivalent to continuity as a function of $x,y$. The triangle
inequality shows that continuity of $f$ implies continuity of $|f|$.
Here attaining a maximum at $z_0$ means $|f(z)|\leq|f(z_0)|$ for all
$z\in E$, with the inequality reversed for a minimum.
\begin{theorem}
A continuous complex function on a compact set is bounded and attains a
maximum and a minimum of its absolute value.
\end{theorem}
\src{44}
\begin{example}
The function $z^2$ is continuous on $\CC$. The principal square root is
continuous on $\CC\setminus(-\infty,0]$. Recall its two values
\[
 z_k=\sqrt{|z|}
 \left(\cos\frac{\theta+2k\pi}{2}
 +i\sin\frac{\theta+2k\pi}{2}\right),\quad k=0,1.
\]
For $z=x+iy$ with fixed $x<0$, as $y\to0^+$ the argument tends to $\pi$
and the principal root tends to $i\sqrt{|x|}$; as $y\to0^-$, the argument
tends to $-\pi$ and the root tends to $-i\sqrt{|x|}$. This exhibits the
discontinuity across the negative real axis.
\end{example}
\editorial{At $x=0$ both limits are $0$; the discontinuity concerns $x<0$.
This specifies the endpoint qualification in the source's $x\leq0$ wording.}
\section{Complex differentiation}
\src{45}
\begin{definition}
For $f:\Omega\to\CC$ on an open set, the derivative at $z_0\in\Omega$ is
\[
 f'(z_0)=\lim_{z\to z_0}\frac{f(z)-f(z_0)}{z-z_0},
\]
if this limit exists. Explicitly, every $\varepsilon>0$ must admit $\delta>0$
such that
\[
 \left|\frac{f(z)-f(z_0)}{z-z_0}-f'(z_0)\right|<\varepsilon
 \qquad(z\in D'(z_0,\delta)).
\]
If this holds at every point, $f$ is holomorphic or analytic on $\Omega$.
The class is denoted by $\Hol(\Omega)$. A function holomorphic on all of
$\CC$ is entire.
\end{definition}
\src{46}
If $f,g\in\Hol(\Omega)$, then $f+g,fg\in\Hol(\Omega)$, so
$\Hol(\Omega)$ is a ring. If $g$ never vanishes, then $f/g$ is holomorphic.
For $\alpha,\beta\in\CC$,
\[
 (\alpha f+\beta g)'=\alpha f'+\beta g',\quad
 (fg)'=f'g+fg',\quad (f/g)'=\frac{f'g-fg'}{g^2}.
\]
Compositions of holomorphic functions are also holomorphic.
\src{47}
\begin{lemma}[Chain rule]
If $f\in\Hol(\Omega)$, $f(\Omega)\subseteq\Omega_1$, and
$g\in\Hol(\Omega_1)$, then $h=g\circ f\in\Hol(\Omega)$ and
$h'(z_0)=g'(f(z_0))f'(z_0)$.
\end{lemma}
\src{48}
\begin{theorem}
Let $g$ be analytic on an open set $\Omega_1$ and $f$ continuous on an
open set $\Omega$. Suppose $f(\Omega)\subseteq\Omega_1$, $g'$ never
vanishes, and $g(f(z))=z$ for every $z\in\Omega$ (so $f$ is one-to-one).
Then $f$ is analytic and $f'=1/(g'\circ f)$.
\end{theorem}
\src{49}
\begin{example}
The function $z$ is holomorphic, with derivative $1$. Every polynomial
\[
 p(z)=a_0+a_1z+\cdots+a_nz^n,\qquad
 p'(z)=a_1+2a_2z+\cdots+na_nz^{n-1},
\]
is entire. The function $1/z$ is holomorphic on every open set not containing
$0$, with derivative $-1/z^2$. By contrast, $f(z)=\bar z$ is not holomorphic:
\[
 \frac{f(z_0+h)-f(z_0)}h=\frac{\bar h}{h},
\]
which has different limits along real and purely imaginary $h\to0$.
\end{example}
\section{Real differentiability and the Cauchy--Riemann equations}
\src{50}
\begin{definition}
For $E\subseteq\RR^n$ open and $F:E\to\RR^m$, differentiability at $x$
means that there is a linear map $A:\RR^n\to\RR^m$ such that
\[
 \lim_{h\to0}\frac{|F(x+h)-F(x)-Ah|}{|h|}=0.
\]
We write $A=F'(x)$. Differentiability in $E$ means differentiability at every
point. Equivalently,
\[
 F(x+h)-F(x)=F'(x)h+r(h),\qquad \lim_{h\to0}\frac{|r(h)|}{|h|}=0.
\]
\end{definition}
\src{51}
For $F=(F_1,\ldots,F_m)$, the derivative matrix, when it exists, is
\[
 F'(x_0)=\begin{bmatrix}
 \frac{\partial F_1}{\partial x_1}(x_0)&\cdots&
 \frac{\partial F_1}{\partial x_n}(x_0)\\
 \vdots&\ddots&\vdots\\
 \frac{\partial F_m}{\partial x_1}(x_0)&\cdots&
 \frac{\partial F_m}{\partial x_n}(x_0)
 \end{bmatrix}.
\]
If all partial derivatives exist near $x_0$ and are continuous at $x_0$,
then $F$ is differentiable there. Existence of partial derivatives alone
does not suffice: the function
\[
 F(x,y)=\begin{cases}y^3/(x^2+y^2),&(x,y)\ne(0,0),\\0,&(x,y)=(0,0)\end{cases}
\]
has both partial derivatives at the origin but is not differentiable
there; its partial derivatives are not continuous there.
\editorial{The slide introduces this example as a failed converse.
The example specifically shows that existence of partial derivatives is
insufficient; it is not an example of a differentiable function with
discontinuous partial derivatives.}
\src{52}
Recall that $\bar z$ is not complex differentiable because its difference
quotient is $\bar h/h$. As a real-variable map it is $F(x,y)=(x,-y)$,
which is indefinitely differentiable. Thus real differentiability does not
guarantee holomorphicity. More generally, to $f=u+iv$ we associate
$F(x,y)=(u(x,y),v(x,y))$.
\src{53}
For $h=(h_1,h_2)\to0$, real differentiability gives
\[
\begin{aligned}
u(x_0+h_1,y_0+h_2)-u(x_0,y_0)
 &=u_x(x_0,y_0)h_1+u_y(x_0,y_0)h_2+r_1(h),\\
v(x_0+h_1,y_0+h_2)-v(x_0,y_0)
 &=v_x(x_0,y_0)h_1+v_y(x_0,y_0)h_2+r_2(h),
\end{aligned}
\]
where $r(h)=(r_1(h),r_2(h))$ and $|r_j(h)|/|h|\to0$ for $j=1,2$.
\src{54}
\begin{theorem}[Cauchy--Riemann equations]
If $f:\Omega\to\CC$ is holomorphic, then
\[
 \frac{\partial f}{\partial x}=\frac1i\frac{\partial f}{\partial y}.
\]
Writing $f=u+iv$ with $u,v$ real valued gives
$u_x=v_y$ and $u_y=-v_x$.
These are the Cauchy--Riemann equations.
\end{theorem}
\src{57}
\begin{theorem}
If $f=u+iv$ on an open $\Omega$, and $u,v$ are real differentiable and
satisfy the Cauchy--Riemann equations throughout $\Omega$, then $f$ is
holomorphic.
\end{theorem}
\src{59}
\begin{example}
The real differentiability hypothesis is essential. Consider
$f(x,y)=\sqrt{|xy|}$ as a complex function with zero imaginary part.
It has both partial derivatives at $(0,0)$ and satisfies the
Cauchy--Riemann equations there. It is not real differentiable at the
origin, however, and is not complex differentiable there.
\end{example}

\chapter{Power Series and Integration over Curves}
\lecturedetails{September 11, 2025}{60--86}
\section{Power series}
\src{61}
Given $(w_n)_{n\geq0}\subseteq\CC$, the series $\sum_{n=0}^\infty w_n$
converges to $w$ if $\lim_{n\to\infty}\sum_{k=0}^nw_k=w$; otherwise it
diverges. It converges if and only if its partial sums are Cauchy, that is,
$\lim_{m,n\to\infty}\sum_{k=m}^nw_k=0$.
It converges \emph{absolutely} if $\sum|w_n|$ converges.
As for real series, absolute convergence implies convergence.
Absolute convergence is equivalent to boundedness of the partial sums
$\sum_{k=0}^n|w_k|$.
\src{62}
\begin{theorem}[Ratio test]
For a series $\sum w_n$ of nonzero terms, if
$\limsup_{n\to\infty}|w_{n+1}/w_n|<1$, it converges absolutely.
If $|w_{n+1}/w_n|\geq1$ for all sufficiently large $n$, it diverges.
\end{theorem}
\begin{exercise}Prove the ratio test.\end{exercise}
\begin{theorem}[Root test]
For any complex series $\sum w_n$, if
$\limsup_{n\to\infty}|w_n|^{1/n}<1$, it converges absolutely.
If this limsup is greater than $1$, it diverges.
\end{theorem}
\begin{exercise}Prove the root test.\end{exercise}
\src{63}
For complex-valued functions $(f_n)$ on a set $S$, pointwise convergence
is equivalent to the pointwise Cauchy condition. Uniform convergence is
equivalent to the uniform Cauchy condition:
$|f_n(z)-f_m(z)|\to0$ uniformly for $z\in S$ as $m,n\to\infty$.
\begin{theorem}[Weierstrass M-test]
If $|g_n(z)|\leq M_n$ on $S\subseteq\CC$ and
$\sum_{n=1}^\infty M_n<\infty$, then $\sum_{n=1}^\infty g_n(z)$
converges uniformly on $S$.
\end{theorem}
\src{64}
Series $\sum_{n=0}^\infty a_n(z-z_0)^n$, with $a_n,z_0\in\CC$, are
called \emph{power series}. They are series of functions of the special
form $f_n(z)=a_n(z-z_0)^n$.
\begin{theorem}
If such a series converges at $z$ with $|z-z_0|=r$, it converges absolutely
in $D(z_0,r)$, uniformly on every closed subdisc and hence on every compact
subset of $D(z_0,r)$.
\end{theorem}
\src{65}
\begin{theorem}[Radius of convergence]
The radius of convergence of $\sum a_n(z-z_0)^n$ is
\[
 r=\left(\limsup_{n\to\infty}\sqrt[n]{|a_n|}\right)^{-1},
\]
with conventions $1/0=\infty$ and $1/\infty=0$. The series converges
absolutely on $D(z_0,r)$, uniformly on compact subsets, and diverges
when $|z-z_0|>r$.
\end{theorem}
\src{66}
\begin{definition}
A function $f:\Omega\to\CC$ on an open set is representable by power
series in $\Omega$ if, for every disc $D(a,r)\subseteq\Omega$, there are
coefficients $c_n$ with
\[
 f(z)=\sum_{n=0}^\infty c_n(z-a)^n\qquad(z\in D(a,r)).
\]
\end{definition}
\begin{theorem}
Such a function is holomorphic, and its derivative is also representable
by power series. On the same disc,
\[
 f'(z)=\sum_{n=1}^\infty nc_n(z-a)^{n-1}.
\]
\end{theorem}
\src{70}
\begin{remark}
The theorem applies again to $f'$, so derivatives of every order exist and
are representable by power series:
\[
 f^{(k)}(z)=\sum_{n=k}^\infty
 n(n-1)\cdots(n-k+1)c_n(z-a)^{n-k}.
\]
Consequently $k!c_k=f^{(k)}(a)$ for $k=0,1,2,\ldots$, which ensures the
uniqueness of the coefficient sequence at each center.
\end{remark}
\src{71}
\begin{example}[Geometric series]
The series $\sum_{n=0}^\infty z^n$ converges absolutely exactly for $|z|<1$.
Its sum there is $1/(1-z)$, although that function is holomorphic on
$\CC\setminus\{1\}$. As in the real case,
\[
 \sum_{n=0}^Nz^n=\frac{1-z^{N+1}}{1-z},\qquad z^{N+1}\to0\quad(|z|<1).
\]
Termwise differentiation gives, on $D(0,1)$,
\[
 \frac1{(1-z)^2}=\left(\frac1{1-z}\right)'
 =\sum_{n=1}^\infty nz^{n-1}.
\]
\end{example}
\src{72--75}
\begin{example}[Exponential and trigonometric functions]
The complex exponential, cosine, and sine are defined by
\[
 e^z=\exp z=\sum_{n=0}^\infty\frac{z^n}{n!},\quad
 \cos z=\sum_{n=0}^\infty(-1)^n\frac{z^{2n}}{(2n)!},\quad
 \sin z=\sum_{n=0}^\infty(-1)^n\frac{z^{2n+1}}{(2n+1)!}.
\]
They agree with the usual real functions on $\RR$, and all three series
converge absolutely on $\CC$. Indeed $|z^n/n!|=|z|^n/n!$, so the
exponential series is bounded in absolute value by
$\sum|z|^n/n!=e^{|z|}<\infty$. This also proves uniform convergence on
each bounded disc. The same argument applies to sine and cosine.
\[
 \exp'(z)=\sum_{n=1}^\infty n\frac{z^{n-1}}{n!}
 =\sum_{m=0}^\infty\frac{z^m}{m!}=\exp z,\qquad
 \cos' z=-\sin z,\quad \sin' z=\cos z.
\]
These functions are entire. Absolute convergence permits multiplication:
\[
\begin{aligned}
\left(\sum_{k=0}^\infty\frac{z^k}{k!}\right)
\left(\sum_{m=0}^\infty\frac{w^m}{m!}\right)
 &=\sum_{n=0}^\infty\frac1{n!}
   \sum_{k=0}^n\frac{n!}{k!(n-k)!}z^kw^{n-k}\\
 &=\sum_{n=0}^\infty\frac{(z+w)^n}{n!}.
\end{aligned}
\]
Thus $\exp z\exp w=\exp(z+w)$, and a direct calculation gives
$\exp(iy)=\cos y+i\sin y$ for real $y$.
Moreover $\exp z=1$ if and only if $z=2k\pi i$, $k\in\ZZ$.
For $z=x+iy$, the modulus condition and the trigonometric identity force
$x=0$; then $\cos y=1$ and $\sin y=0$ force $y=2k\pi$.
Consequently the complex exponential has period $2\pi i$:
$\exp(z+2k\pi i)=\exp z$ for $z\in\CC$, $k\in\ZZ$.
\end{example}
\section{Integration over paths}
\src{76}
\begin{definition}[Curve]
In a topological space $X$, a curve is a continuous map
$\gamma:[\alpha,\beta]\to X$, where $\alpha<\beta$ and the interval is
compact. Its parameter interval is $[\alpha,\beta]$ and its range is
$\gamma^*=\gamma([\alpha,\beta])=\{\gamma(t):t\in[\alpha,\beta]\}$.
\end{definition}
The range is compact and connected. The curve is closed if its initial
point $\gamma(\alpha)$ equals its endpoint $\gamma(\beta)$.
We distinguish the oriented geometric object $\gamma^*$ from its
parametrization $\gamma$; the parametrization is not unique.
\src{77}
Two parametrizations $\gamma:[\alpha,\beta]\to\CC$ and
$\widetilde\gamma:[\alpha_1,\beta_1]\to\CC$ are equivalent if there is a
continuously differentiable bijection
$\phi:[\alpha_1,\beta_1]\to[\alpha,\beta]$ with $\phi'>0$ and
$\widetilde\gamma(s)=\gamma(\phi(s))$.
The positive derivative preserves orientation as $s$ travels from
$\alpha_1$ to $\beta_1$. Equivalent parametrizations give the same oriented
curve, whose orientation is determined by increasing parameter.
\src{78}
\begin{definition}[Path]
A path is a piecewise continuously differentiable plane curve. More precisely,
it is a continuous complex function $\gamma$ on $[\alpha,\beta]$ for which
there are finitely many points
$\alpha=s_0<s_1<\cdots<s_n=\beta$, and $\gamma$ has a continuous derivative
on each $[s_{j-1},s_j]$. Left and right derivatives at the intervening
endpoints may differ. A closed path is a closed curve which is a path.
\end{definition}
\src{79}
\begin{definition}[Path integral]
If $f$ is continuous on the trace of a path $\gamma$, define
\[
 \int_\gamma f(z)\dd z=\int_\alpha^\beta f(\gamma(t))\gamma'(t)\dd t.
\]
The integral on the right is a Riemann integral: $\gamma'$ is bounded
and has at most finitely many discontinuities.
\end{definition}
Unless otherwise mentioned, closed contour integrals in the lectures use
the anticlockwise orientation.
\src{80}
Let $\phi:[a,b]\to[\alpha,\beta]$ be continuous, strictly increasing,
onto, and continuously differentiable. Then
$\phi(a)=\alpha$, $\phi(b)=\beta$, and $\sigma=\gamma\circ\phi$ is a
path on $[a,b]$. For $f$ continuous on the common trace,
$\int_\sigma f(z)\dd z=\int_\gamma f(z)\dd z$.
The map $\phi$ is a change of parameter.
\src{81}
If the endpoint of $\gamma_1$ is the initial point of $\gamma_2$, a suitable
reparametrization allows us to follow first $\gamma_1$, then $\gamma_2$.
The resulting path is denoted $\gamma=\gamma_1+\gamma_2$, and
\[
 \int_\gamma f=\int_{\gamma_1}f+\int_{\gamma_2}f
\]
for continuous $f$ on $\gamma_1^*\cup\gamma_2^*$.
For paths $\gamma_1,\ldots,\gamma_n$ with matching successive endpoints,
$\gamma_1+\cdots+\gamma_n$ is defined similarly.
\src{82}
For $\gamma:[0,1]\to\CC$, the opposite path is
$\gamma_1(t)=\gamma(1-t)$ and $\int_\gamma f=-\int_{\gamma_1}f$.
Also, for a path on $[\alpha,\beta]$,
\[
\begin{aligned}
 \left|\int_\gamma f(z)\dd z\right|
 &=\left|\int_\alpha^\beta f(\gamma(t))\gamma'(t)\dd t\right|\\
 &\leq\max_{z\in\gamma^*}|f(z)|\int_\alpha^\beta|\gamma'(t)|\dd t
 =\ell(\gamma)\max_{z\in\gamma^*}|f(z)|.
\end{aligned}
\]
Here $\ell(\gamma)=\int_\alpha^\beta|\gamma'(t)|\dd t$ is the length.
\src{83}
\begin{remark}
For a closed path, $\Chat\setminus\gamma^*$ is open and is a disjoint
union of regions. These are the regions determined by $\gamma$.
Precisely one contains $\infty$, corresponding to the unbounded region.
The regions determined in $\Chat$ and in $\CC$ are identical except that
the unbounded region in $\CC$ does not contain $\infty$.
\end{remark}
\src{84}
\begin{example}[Circle]
The positively oriented circle of center $a$ and radius $r$ is
$\gamma(t)=a+re^{it}$, $0\leq t\leq2\pi$. Its length is $2\pi r$ and
\[
 \int_\gamma f(z)\dd z=\int_0^{2\pi}f(a+re^{it})ire^{it}\dd t.
\]
\end{example}
\src{85}
\begin{example}[Segment]
The positively oriented interval $[a,b]$ is parametrized by
$\gamma(t)=a+(b-a)t$, $0\leq t\leq1$. It has length $|b-a|$ and
\[
 \int_{[a,b]}f(z)\dd z=(b-a)\int_0^1f(a+(b-a)t)\dd t.
\]
On an arbitrary $[\alpha,\beta]$, the equivalent parametrization is
$\gamma_1(t)=[a(\beta-t)+b(t-\alpha)]/(\beta-\alpha)$.
The opposite path is $[b,a]$.
\end{example}
\src{86}
\begin{example}[Triangle]
For an ordered triple $\{a,b,c\}$, the triangle $\Delta(a,b,c)$ is the
smallest convex set containing these points. For continuous $f$ on its
boundary, define
\[
 \int_{\partial\Delta}f=\int_{[a,b]}f+\int_{[b,c]}f+\int_{[c,a]}f.
\]
Alternatively, regard $\partial\Delta$ as the concatenation of these paths,
and the equality follows from path addition. Cyclic permutation of
$a,b,c$ does not change the integral; replacing the order by $a,c,b$
changes its sign.
\end{example}

\chapter{Index function and Cauchy's theorem}
\lecturedetails{September 15, 2025}{87--113}
\section{Index and primitives}
\src{88}
\begin{definition}
For a path $\gamma$ and $a\notin\gamma^*$, write
\[
\Ind_\gamma(a)=\frac1{2\pi i}\int_\gamma\frac{dz}{z-a}.
\]
This is the \emph{index}, or \emph{winding number}, of $\gamma$ about $a$.
\end{definition}
\begin{theorem}
If $\gamma$ is closed and $\Omega=\CC\setminus\gamma^*$, then
$\Ind_\gamma$ is integer-valued, constant on each region of $\Omega$,
and zero on the unbounded region.
\end{theorem}
\src{91}
\begin{remark}
For a closed curve and $a$ outside its image, the number of times
the curve winds around $a$ can be computed from the change of argument
of $z-a$ as $z$ travels along the curve. The quantity $\arg(z-a)/(2\pi)$
increases or decreases by one for each complete positive or negative turn.
Logarithms will later give a geometric explanation of this interpretation.
\end{remark}
\src{92}
\begin{theorem}
For the positively oriented circle $\gamma$ of center $a$ and radius $r$,
\[
\Ind_\gamma(z)=
\begin{cases}1,&|z-a|<r,\\0,&|z-a|>r.\end{cases}
\]
\end{theorem}
\src{93}
\begin{definition}
A \emph{primitive} of $f:\Omega\to\CC$, where $\Omega$ is open,
is a holomorphic function $F$ on $\Omega$ with $F'=f$.
\end{definition}
\begin{theorem}
If $f$ is continuous on $\Omega$ and has a primitive $F$ there,
then for any path $\gamma$ from $w_1$ to $w_2$ in $\Omega$,
$\int_\gamma f(z)\,dz=F(w_2)-F(w_1)$.
\end{theorem}
\src{95}
\begin{corollary}
If $f$ is continuous on an open set, has a primitive there, and
$\gamma$ is a closed path in that set, then $\int_\gamma f(z)\,dz=0$.
\end{corollary}
\begin{example}
The function $f(z)=1/z$ has no primitive on $\CC\setminus\{0\}$.
On the unit circle $C$, parametrized by $z=e^{it}$,
\[
\int_C f(z)\,dz=\int_0^{2\pi}\frac{ie^{it}}{e^{it}}\,dt=2\pi i\ne0.
\]
\end{example}
\src{96}
\begin{corollary}
If $f$ is holomorphic on a region $\Omega$ and $f'=0$, then $f$ is constant.
\end{corollary}
\section{Local Cauchy theorem}
\src{97}
\begin{theorem}[Cauchy--Goursat]
Let $\Omega\subseteq\CC$ be open, let $\Delta$ be a closed triangle
contained in $\Omega$, and let $p\in\Omega$. If $f$ is continuous
on $\Omega$ and holomorphic on $\Omega\setminus\{p\}$, then
$\int_{\partial\Delta}f(z)\,dz=0$.
\end{theorem}
\begin{remark}
In particular, this holds for every triangular path in $\Omega$
when $f$ is holomorphic on $\Omega$. This is Cauchy's theorem
for triangular paths.
\end{remark}
\src{104}
\begin{theorem}[Cauchy's theorem for convex sets]
Let $\Omega\subseteq\CC$ be convex and open, and let $p\in\Omega$.
If $f$ is continuous on $\Omega$ and holomorphic on
$\Omega\setminus\{p\}$, then $f$ has a primitive on $\Omega$, and
$\int_\gamma f(z)\,dz=0$ for every closed path $\gamma$ in $\Omega$.
\end{theorem}
\begin{corollary}
Every holomorphic function on an open convex set has a primitive;
its integral along every closed path in that set is zero.
\end{corollary}
\src{107}
\begin{remark}
Cauchy's theorem also holds for the illustrated regions,
although some of them are not convex.
\end{remark}
\sourcefigure{p107}{.93}{The original multiple keyhole, rectangular keyhole,
semicircle, indented semicircle, sector, and parallelogram (slide 107).}
\begin{exercise}
Explain why Cauchy's theorem holds in each of these regions.
\end{exercise}
\section{Worked contour-integral examples}
\src{108--110}
\begin{example}[Gaussian Fourier transform]
For $\xi\in\RR$,
\[
e^{-\pi\xi^2}=\int_{-\infty}^{\infty}e^{-\pi x^2}e^{-2\pi ix\xi}\,dx.
\]
For $\xi=0$, this is the known Gaussian integral
$\int_\RR e^{-\pi x^2}\,dx=1$. For $\xi>0$, use the entire
function $f(z)=e^{-\pi z^2}$ and the positively oriented rectangle
with vertices $R,R+i\xi,-R+i\xi,-R$.
\sourcefigure{p108}{.93}{The rectangle $\gamma_R$ used for the Gaussian transform.}
Cauchy's theorem gives $\int_{\gamma_R}f(z)\,dz=0$.
The bottom side contributes $\int_{-R}^R e^{-\pi x^2}\,dx\to1$.
On the right side,
\[
I(R)=\int_0^\xi f(R+iy)i\,dy
=i\int_0^\xi e^{-\pi(R^2+2iRy-y^2)}\,dy.
\]
For fixed $\xi$, $|I(R)|\le C e^{-\pi R^2}\to0$.
The left side has the same vanishing limit. The top contributes
\[
\int_R^{-R}e^{-\pi(x+i\xi)^2}\,dx
=-e^{\pi\xi^2}\int_{-R}^R e^{-\pi x^2}e^{-2\pi ix\xi}\,dx.
\]
Thus the limiting identity is
$0=1-e^{\pi\xi^2}\int_\RR e^{-\pi x^2}e^{-2\pi ix\xi}\,dx$.
For $\xi<0$, use the corresponding rectangle in the lower half-plane.
\end{example}
\src{111--113}
\begin{example}
Another classical integral is
$\int_0^\infty(1-\cos x)/x^2\,dx=\pi/2$.
Take $f(z)=(1-e^{iz})/z^2$ and an indented semicircle
in the upper half-plane.
\sourcefigure{p111}{.93}{The small negatively oriented and large positively
oriented semicircles in the original example.}
Denote the small and large arcs by $\gamma_\varepsilon^+$ and
$\gamma_R^+$; their orientations are negative and positive, respectively.
Cauchy's theorem yields
\[
\int_{-R}^{-\varepsilon}f(x)\,dx+
\int_{\gamma_\varepsilon^+}f(z)\,dz+
\int_\varepsilon^R f(x)\,dx+
\int_{\gamma_R^+}f(z)\,dz=0.
\]
On the large arc, $|f(z)|\le2/|z|^2$, so its integral tends to zero
as $R\to\infty$. Consequently
\[
\int_{|x|\ge\varepsilon}\frac{1-e^{ix}}{x^2}\,dx
=-\int_{\gamma_\varepsilon^+}\frac{1-e^{iz}}{z^2}\,dz.
\]
Near zero, $f(z)=-i/z+E(z)$, where $E$ is bounded.
Parametrizing the small arc by $z=\varepsilon e^{i\theta}$,
with $\theta$ decreasing from $\pi$ to $0$, gives
\[
\lim_{\varepsilon\to0}
\int_{\gamma_\varepsilon^+}f(z)\,dz
=\int_\pi^0(-i)i\,d\theta=-\pi.
\]
Taking real parts gives $\int_\RR(1-\cos x)/x^2\,dx=\pi$.
The integrand is even, giving the stated half-line integral.
\end{example}

\chapter{Cauchy's integral formula and Laurent series}
\lecturedetails{September 18, 2025}{114--139}
\section{Integral formula and its consequences}
\src{115}
\begin{theorem}[Recall: Cauchy's theorem for convex sets]
If $\Omega$ is open and convex, $p\in\Omega$, and $f$ is continuous
on $\Omega$ and holomorphic on $\Omega\setminus\{p\}$, then $f$
has a primitive on $\Omega$ and $\int_\gamma f(z)\,dz=0$ for
every closed path $\gamma$ in $\Omega$.
\end{theorem}
\begin{theorem}[Cauchy's integral formula for convex sets]
Let $\gamma$ be a closed path in an open convex set $\Omega\subseteq\CC$
and $f\in H(\Omega)$. For every $a\in\Omega\setminus\gamma^*$,
\[
\frac1{2\pi i}\int_\gamma\frac{f(z)}{z-a}\,dz
=\Ind_\gamma(a)f(a).
\]
\end{theorem}
\src{117}
\begin{corollary}
If $f\in H(\Omega)$, then $f^{(n)}\in H(\Omega)$ for every $n\in\NN$.
If a positively oriented circle $C$ and its interior lie in $\Omega$,
then, for $z$ in that interior,
\[
f^{(n)}(z)=\frac{n!}{2\pi i}\int_C\frac{f(\zeta)}{(\zeta-z)^{n+1}}\,d\zeta.
\]
Thus holomorphic functions have infinitely many complex derivatives.
\end{corollary}
\src{120}
\begin{corollary}[Cauchy's inequality]
If $f$ is holomorphic on an open set containing $\overline{D(z_0,R)}$,
and $C=\partial D(z_0,R)$, then
\[
|f^{(n)}(z_0)|\le\frac{n!}{R^n}\|f\|_{L^\infty(C)},\qquad
\|f\|_{L^\infty(C)}=\sup_{z\in C}|f(z)|.
\]
\end{corollary}
\src{121}
\begin{theorem}
If $f\in H(\Omega)$ and $D$ is a disc centered at $z_0$ with
$\overline D\subset\Omega$, then, for every $z\in D$,
\[
f(z)=\sum_{n=0}^{\infty}a_n(z-z_0)^n,\qquad
a_n=\frac{f^{(n)}(z_0)}{n!}\quad(n\ge0).
\]
\end{theorem}
Since power series are indefinitely complex differentiable, this
provides another route to the infinite differentiability of holomorphic functions.
\src{124}
\begin{remark}
The expansion about $z_0$ converges in any disc, however large,
whose closure lies in $\Omega$. In particular, an entire function
has an expansion $f(z)=\sum_{n=0}^\infty a_nz^n$ converging on all of $\CC$.
\end{remark}
\src{125}
\begin{theorem}[Concluding equivalences]
For an open set $\Omega\subseteq\CC$, the following are equivalent:
\begin{enumerate}
\item $f\in H(\Omega)$.
\item $f^{(n)}\in H(\Omega)$ for every $n\ge0$.
\item On every disc $D$ centered at $z_0$ with $\overline D\subset\Omega$,
$f(z)=\sum_{n\ge0}a_n(z-z_0)^n$, with $a_n=f^{(n)}(z_0)/n!$.
\end{enumerate}
\end{theorem}
\section{Laurent series}
\src{126}
\begin{definition}
For $0\le s_1<s_2\le+\infty$,
\[
A(z_0,s_1,s_2)=\{z:s_1<|z-z_0|<s_2\}
\]
is an open annulus. For finite radii, the closed annulus is
$\overline A(z_0,s_1,s_2)=\{z:s_1\le|z-z_0|\le s_2\}$.
In the usual positive finite-radius case,
\begin{align*}
A(z_0,s_1,s_2)&=D(z_0,s_2)\setminus\overline{D(z_0,s_1)},\\
\overline A(z_0,s_1,s_2)&=\overline{D(z_0,s_2)}\setminus D(z_0,s_1),\\
\overline A(z_0,s_1,s_2)&=A(z_0,s_1,s_2)\cup C(z_0,s_1)\cup C(z_0,s_2).
\end{align*}
\end{definition}
\editorial{At radius zero, use the distance definitions;
the inner boundary degenerates to the center. An infinite radius
does not give a finite boundary circle.}
\src{127}
\begin{definition}
A \emph{Laurent series} about $z_0$ is
$\sum_{n\in\ZZ}a_n(z-z_0)^n$. Its convergence means that both
$\sum_{n\ge0}a_n(z-z_0)^n$ and $\sum_{n<0}a_n(z-z_0)^n$
converge; its value is their sum. It diverges if either tail diverges.
\end{definition}
\src{128}
\begin{theorem}
For a Laurent series, put
\[
r=\limsup_{n\to\infty}|a_{-n}|^{1/n},\qquad
R=\left(\limsup_{n\to\infty}|a_n|^{1/n}\right)^{-1}.
\]
If $r<R$, the series converges absolutely on $A(z_0,r,R)$,
uniformly on its compact subsets, and defines a holomorphic function there.
\end{theorem}
\src{131}
\begin{theorem}[Cauchy's formula for annuli]
Let $0<r_1<r_2<\infty$, and let $f$ be holomorphic on an open
neighborhood of $\overline A(z_0,r_1,r_2)$. If $\gamma_1,\gamma_2$
are its positively oriented inner and outer circles, then
\[
f(z)=\frac1{2\pi i}\int_{\gamma_2}\frac{f(w)}{w-z}\,dw
-\frac1{2\pi i}\int_{\gamma_1}\frac{f(w)}{w-z}\,dw
\quad(z\in A(z_0,r_1,r_2)).
\]
\end{theorem}
\begin{exercise}[Question retained from the omitted argument]
Explain in what sense Cauchy's integral formula is available
on the annular region $A(z_0,s_1,s_2)$ with
$0<s_1<r_1<r_2<s_2$.
\end{exercise}
\editorial{Use the annular boundary formula, for example by cutting
along a radial segment. The assertion that every holomorphic function
on an annulus has a primitive would be false.}
\src{133}
\begin{theorem}[Laurent representation]
If $f\in H(A(z_0,s_1,s_2))$, there is a unique sequence
$(a_n)_{n\in\ZZ}$ such that
\[
f(z)=\sum_{n=-\infty}^{\infty}a_n(z-z_0)^n.
\]
For $s_1<r<s_2$ and $n\in\ZZ$,
\[
a_n=\frac1{2\pi i}\int_{C(z_0,r)}
\frac{f(w)}{(w-z_0)^{n+1}}\,dw.
\]
The expansion converges absolutely on the annulus and uniformly
on its compact subsets.
\end{theorem}
\src{138}
\begin{exercise}[Question retained from the omitted argument]
Suppose $f(z)=\sum_{n\in\ZZ}c_n(z-z_0)^n$ for every
$z\in A(z_0,s_1,s_2)$. Why must this series converge uniformly
on compact subsets of that annulus?
\end{exercise}

\chapter{Further applications and singularities}
\lecturedetails{September 22, 2025}{140--165}
\section{Further applications of Cauchy's formula}
\src{141}
\begin{corollary}[Cauchy's inequality, recalled]
If $f$ is holomorphic on an open neighborhood of $\overline{D(z_0,R)}$
and $C=\partial D(z_0,R)$, then
\[
|f^{(n)}(z_0)|\le n!R^{-n}\|f\|_{L^\infty(C)},\qquad
\|f\|_{L^\infty(C)}=\sup_{z\in C}|f(z)|.
\]
\end{corollary}
\src{142}
\begin{theorem}[Liouville]
Every bounded entire function is constant.
\end{theorem}
\begin{theorem}[Fundamental theorem of algebra]
Every nonconstant polynomial $P(z)=a_nz^n+\cdots+a_1z+a_0$
with complex coefficients has a root in $\CC$. If its degree is
$n\ge1$, it has exactly $n$ roots, counted with multiplicity.
\end{theorem}
\src{144}
\begin{theorem}[Morera: converse to Cauchy's theorem]
If $f$ is continuous on an open set $\Omega\subseteq\CC$ and
$\int_{\partial\Delta}f(z)\,dz=0$ for every closed triangle
$\Delta\subset\Omega$, then $f\in H(\Omega)$.
\end{theorem}
\src{146}
\begin{definition}
The sequence $(f_j)$ converges to $f$ \emph{uniformly on compact subsets}
of $\Omega$ if for every compact $K\subseteq\Omega$ and $\varepsilon>0$
there exists $N=N(K,\varepsilon)$ such that
$\sup_{z\in K}|f_j(z)-f(z)|<\varepsilon$ whenever $j>N$.
\end{definition}
\begin{theorem}
If $(f_j)\subset H(\Omega)$ converges uniformly on compact subsets to $f$,
then $f\in H(\Omega)$ and $f_j'\to f'$ uniformly on compact subsets.
\end{theorem}
\section{Zeros and analytic continuation}
\src{148}
\begin{theorem}
Let $\Omega$ be a region, $f\in H(\Omega)$, and
$Z(f)=\{a\in\Omega:f(a)=0\}$. Either $Z(f)=\Omega$, or $Z(f)$ has
no limit point in $\Omega$. In the latter case, for each $a\in Z(f)$
there is a unique positive integer $m=m(a)$ with
\[
f(z)=(z-a)^m g(z)\quad(z\in\Omega),\qquad
g\in H(\Omega),\quad g(a)\ne0.
\]
Moreover, $Z(f)$ is at most countable. The integer $m$ is the
\emph{order} of the zero at $a$.
\end{theorem}
\src{153}
\begin{theorem}[Identity theorem]
If $f,g$ are holomorphic on a region $\Omega$ and agree on a set
having a limit point in $\Omega$, then $f=g$ throughout $\Omega$.
\end{theorem}
\begin{definition}
If $f,F$ are analytic on regions $\Omega,\Omega'$ with
$\Omega\subseteq\Omega'$ and $F=f$ on $\Omega$, then $F$ is an
\emph{analytic continuation} of $f$ to $\Omega'$.
The identity theorem guarantees uniqueness of this continuation.
\end{definition}
\src{154--155}
Let $\Omega$ be open and symmetric about the real axis:
$z\in\Omega$ if and only if $\bar z\in\Omega$.
Let $\Omega^+$ and $\Omega^-$ be its upper and lower half-plane
parts, and set $I=\Omega\cap\RR$. Thus
$\Omega=\Omega^+\cup I\cup\Omega^-$. The set $I$ is the relative
interior of the real-axis portion of their common boundary.
\sourcefigure{p154}{.8}{The symmetric region from slide 154.}
\begin{theorem}[Symmetry principle]
If $f^+,f^-$ are holomorphic on $\Omega^+,\Omega^-$, extend continuously
to $I$, and $f^+(x)=f^-(x)$ there, then
\[
f(z)=\begin{cases}
f^+(z),&z\in\Omega^+,\\
f^+(z)=f^-(z),&z\in I,\\
f^-(z),&z\in\Omega^-
\end{cases}
\]
is holomorphic on all of $\Omega$.
\end{theorem}
\src{158}
\begin{theorem}[Schwarz reflection]
If $f$ is holomorphic on $\Omega^+$, extends continuously to $I$,
and is real-valued on $I$, then a holomorphic function $F$ on
$\Omega$ exists with $F=f$ on $\Omega^+$.
\end{theorem}
\section{Classification of singularities}
\src{159}
\begin{definition}
If $a\in\Omega$ and $f\in H(\Omega\setminus\{a\})$, then $a$ is an
\emph{isolated singularity}. It is \emph{removable} if $f$ can be
defined at $a$ so as to be holomorphic on $\Omega$.
\end{definition}
\begin{theorem}[Riemann's removable-singularity theorem]
If $f\in H(\Omega\setminus\{a\})$ is bounded on a punctured disc
$D'(a,r)$, then its singularity at $a$ is removable.
\end{theorem}
\src{161}
\begin{theorem}
For $a\in\Omega$ and $f\in H(\Omega\setminus\{a\})$, exactly one of
the following occurs:
\begin{enumerate}
\item The singularity is removable.
\item There are $c_1,\ldots,c_m\in\CC$, with $m\ge1$ and $c_m\ne0$,
such that $f(z)-\sum_{k=1}^m c_k(z-a)^{-k}$ has a removable singularity.
\item For every $r>0$ with $D(a,r)\subseteq\Omega$,
$f(D'(a,r))$ is dense in $\CC$.
\end{enumerate}
The last conclusion is the Casorati--Weierstrass theorem.
\end{theorem}
\src{162}
\begin{remark}
In the second case, $f$ has a \emph{pole of order $m$}.
The polynomial in $(z-a)^{-1}$, $\sum_{k=1}^m c_k(z-a)^{-k}$,
is its \emph{principal part}; then $|f(z)|\to\infty$ as $z\to a$.
In the third case, the singularity is \emph{essential}.
Equivalently, for every $w\in\CC$ there is a sequence
$z_n\to a$ with $f(z_n)\to w$.
\end{remark}
\src{165}
\begin{remark}
For the Laurent expansion $f(z)=\sum_{n=-\infty}^\infty a_n(z-z_0)^n$:
\begin{itemize}
\item The singularity is removable exactly when $a_n=0$ for every $n<0$.
\item It is a pole of order $m$ when $m$ is the largest positive
integer with $a_{-m}\ne0$. A pole of order one is \emph{simple}.
\item It is essential when infinitely many negative-index coefficients
are nonzero.
\item The behavior at infinity is studied using $g(z)=f(1/z)$ near zero.
The point $\infty$ is an isolated singularity if $f$ is analytic for
$|z|>r$ for some $r$; its type is the type of $g$ at zero.
\end{itemize}
\end{remark}

\chapter{Maximum principle, open mappings, and logarithms}
\lecturedetails{September 25, 2025}{166--195}
\section{Maximum modulus and open mappings}
\src{167}
\begin{theorem}[Orthogonality relation]
If $f(z)=\sum_{n\ge0}c_n(z-a)^n$ on $D(a,R)$ and $0<r<R$, then
\[
\sum_{n\ge0}|c_n|^2r^{2n}
=\frac1{2\pi}\int_{-\pi}^{\pi}|f(a+re^{i\theta})|^2\,d\theta.
\]
\end{theorem}
\src{169}
\begin{definition}
The modulus $|f|$ has a local maximum at $a\in\Omega$ if some
$\delta>0$ satisfies $D(a,\delta)\subseteq\Omega$ and
$|f(a)|\ge|f(z)|$ throughout this disc.
Having no local maximum means having none at any point of $\Omega$.
A local minimum is defined with the inequality reversed.
\end{definition}
\begin{theorem}[Maximum modulus principle]
If $\Omega$ is a region and $f\in H(\Omega)$, then $|f|$ has no
local maximum unless $f$ is constant. If additionally $\overline\Omega$
is compact and $f$ is continuous on $\overline\Omega$, then
$\sup_\Omega|f|\le\sup_{\partial\Omega}|f|$.
\end{theorem}
\src{172}
\begin{exercise}[Corollary to prove]
Let $\Omega$ be a region, $f\in H(\Omega)$, and
$\overline{D(a,r)}\subset\Omega$. Prove that
\[
|f(a)|\le\max_{-\pi\le\theta\le\pi}|f(a+re^{i\theta})|,
\]
with equality if and only if $f$ is constant on $\Omega$.
Deduce the absence of local maxima for nonconstant $|f|$.
If $f$ has no zero in $D(a,r)$, prove also that
\[
|f(a)|\ge\min_{-\pi\le\theta\le\pi}|f(a+re^{i\theta})|.
\]
\end{exercise}
\src{173}
\begin{lemma}
If $f\in H(\Omega)$, then
\[
g(z,w)=\begin{cases}
(f(z)-f(w))/(z-w),&z\ne w,\\f'(z),&z=w
\end{cases}
\]
is continuous on $\Omega\times\Omega$.
\end{lemma}
\src{175}
\begin{theorem}
If $\varphi\in H(\Omega)$ and $\varphi'(z_0)\ne0$, there is a
neighborhood $V\subseteq\Omega$ of $z_0$ on which $\varphi$ is
one-to-one, $W=\varphi(V)$ is open, and the inverse
$\psi:W\to V$, defined by $\psi(\varphi(z))=z$, is holomorphic.
\end{theorem}
\src{178}
\begin{definition}
For $m\in\NN$, write $\pi_m(z)=z^m$.
\end{definition}
Every $w=re^{i\theta}\ne0$ has exactly $m$ distinct preimages:
\[
\pi_m(z)=w\quad\Longleftrightarrow\quad
z=r^{1/m}e^{i(\theta+2k\pi)/m},\quad k=1,\ldots,m.
\]
Each $\pi_m$ is open: away from zero use the preceding theorem,
and at zero use $\pi_m(D(0,r))=D(0,r^m)$.
Compositions of open maps are open. In particular, $\pi_m\circ\varphi$
is open when $\varphi'$ has no zero.
\src{179}
\begin{theorem}[Local form and open mapping theorem]
Let $f$ be nonconstant and holomorphic on a region $\Omega$,
$z_0\in\Omega$, and $w_0=f(z_0)$. If $m$ is the order of the
zero of $f-w_0$ at $z_0$, there are a neighborhood $V$ of $z_0$
and $\varphi\in H(V)$ such that
\[
f(z)=w_0+\varphi(z)^m.
\]
Moreover, $\varphi'$ has no zero on $V$, and $\varphi$ maps
$V$ invertibly onto a disc $D(0,r)$.
\end{theorem}
\begin{remark}
Thus $f-w_0=\pi_m\circ\varphi$ locally.
The map $f:V\setminus\{z_0\}\to D'(w_0,r^m)$ is exactly $m$-to-one.
Every $w_0\in f(\Omega)$ is an interior point, so $f(\Omega)$ is open.
\end{remark}
\src{181}
\begin{theorem}[Inverse mapping theorem]
If $f\in H(\Omega)$ is one-to-one on a region $\Omega$, then
$f'(z)\ne0$ everywhere and $f^{-1}$ is holomorphic.
\end{theorem}
\begin{remark}
The converse is false: $f(z)=e^z$ has nonvanishing derivative
everywhere but is not one-to-one on $\CC$.
\end{remark}
\section{Logarithms and continuous arguments}
\src{182}
The exponential restricted to $S_\alpha=\{x+iy:\alpha\le y<\alpha+2\pi\}$
is a bijection onto $\CC\setminus\{0\}$.
\begin{definition}
Its inverse is $\log_\alpha$, and $\arg_\alpha=\Im\log_\alpha$.
Thus $\log_\alpha(e^z)=z$ on $S_\alpha$ and
$e^{\log_\alpha z}=z$ for $z\ne0$.
The principal branches used here are
$\operatorname{Log}=\log_{-\pi}$ and $\operatorname{Arg}=\arg_{-\pi}$.
\end{definition}
\src{183}
\begin{theorem}
\begin{enumerate}
\item For $z\ne0$, $\log_\alpha z=\log|z|+i\arg_\alpha z$,
where $\arg_\alpha z$ is the unique argument in $[\alpha,\alpha+2\pi)$;
equivalently $z/|z|=e^{i\arg_\alpha z}$.
\item Set $R_\alpha=\{re^{i\alpha}:r\ge0\}$. The two functions are
continuous on $\CC\setminus R_\alpha$ and discontinuous on the
nonzero part of the slit.
\item $\log_\alpha$ is analytic on $\CC\setminus R_\alpha$ and
$(\log_\alpha)'(z)=1/z$.
\end{enumerate}
\end{theorem}
\editorial{Neither function is defined at zero. The half-closed strip
is not an open domain of holomorphy; the exponential is the restriction
of an entire function.}
\src{185}
\begin{theorem}[Recalled inverse criterion]
Let $g$ be analytic on $\Omega_1$ and $f$ continuous on the open set
$\Omega$. If $f(\Omega)\subseteq\Omega_1$, $g'$ is nowhere zero,
and $g(f(z))=z$, then $f$ is analytic and
$f'=1/(g'\circ f)$. The identity also makes $f$ one-to-one.
\end{theorem}
\src{186}
\begin{definition}
Let $S\subseteq\CC$, or more generally let $S$ be a metric space,
and let $f:S\to\CC\setminus\{0\}$ be continuous.
A continuous logarithm is a continuous $g:S\to\CC$ with $f=e^g$.
A continuous argument is a continuous $\theta:S\to\RR$ with
$f=|f|e^{i\theta}$.
\end{definition}
\begin{example}
\begin{enumerate}
\item On $S=[0,2\pi]$, $f(s)=e^{is}$ has continuous arguments
$\theta(s)=s+2k\pi$ for any fixed integer $k$.
\item If $f:S\to\CC\setminus R_\alpha$ is continuous,
$\theta=\arg_\alpha\circ f$ is a continuous argument.
\item On $S=\{z:|z|=1\}$, $f(z)=z$ has no continuous argument.
\end{enumerate}
\end{example}
\src{187}
\begin{theorem}
For continuous nonvanishing $f:S\to\CC$:
\begin{enumerate}
\item If $g$ is a continuous logarithm, $\Im g$ is a continuous argument.
\item If $\theta$ is a continuous argument, $\log|f|+i\theta$ is a
continuous logarithm. Thus either exists if and only if the other does.
\item On connected $S$, any two continuous logarithms differ by
$2\pi ik$ for one fixed $k\in\ZZ$; any two continuous arguments
differ by $2\pi l$ for one fixed $l\in\ZZ$.
\item If $S$ is connected and $s,t\in S$, any continuous logarithm $g$
and argument $\theta$ satisfy
\[
g(s)-g(t)=\log|f(s)|-\log|f(t)|+i(\theta(s)-\theta(t)).
\]
\end{enumerate}
\end{theorem}
\src{189}
\begin{theorem}
Every continuous curve $\gamma:[a,b]\to\CC\setminus\{0\}$
has a continuous argument and hence a continuous logarithm.
\end{theorem}
\src{190}
\begin{theorem}
For a closed curve $\gamma:[a,b]\to\CC$, $z_0\notin\gamma^*$,
and any continuous argument $\theta$ of $\gamma-z_0$,
$\theta(b)-\theta(a)$ is an integer multiple of $2\pi$ and
is independent of the choice of $\theta$.
\end{theorem}
\src{191}
\begin{definition}
The winding number is $W(\gamma,z_0)=(\theta(b)-\theta(a))/(2\pi)$.
\end{definition}
\begin{remark}
It is well-defined by the preceding theorem. Intuitively it is the
net number of revolutions about $z_0$ as the parameter runs from
$a$ to $b$. It is translation-invariant:
$W(\gamma,z_0)=W(\gamma+w,z_0+w)$ for every $w\in\CC$.
\end{remark}
\src{192}
\begin{theorem}
For a closed path,
\[
W(\gamma,z_0)=\frac1{2\pi i}\int_\gamma\frac{dz}{z-z_0}
=\Ind_\gamma(z_0).
\]
More generally, if $f$ is analytic on a neighborhood of $\gamma^*$
and $z_0\notin(f\circ\gamma)^*$, then
\[
W(f\circ\gamma,z_0)=\frac1{2\pi i}
\int_\gamma\frac{f'(z)}{f(z)-z_0}\,dz.
\]
\end{theorem}

\chapter{Global Cauchy's theorem}
\lecturedetails{October 2, 2025}{196--225}
\section{Chains and cycles}
\src{197--198}
For paths $\gamma_1,\ldots,\gamma_n$, put $K=\bigcup_i\gamma_i^*$.
On the vector space $C(K)$ of continuous functions on $K$, each
path induces the linear functional
$\widetilde\gamma_i(f)=\int_{\gamma_i}f(z)\,dz$.
Define $\widetilde\Gamma=\sum_i\widetilde\gamma_i$; explicitly
$\widetilde\Gamma(f)=\sum_i\widetilde\gamma_i(f)$.
This motivates the formal sum $\Gamma=\gamma_1\dotplus\cdots\dotplus\gamma_n$
and the definition
\[
\int_\Gamma f(z)\,dz=\widetilde\Gamma(f)=\sum_{i=1}^n\int_{\gamma_i}f(z)\,dz.
\]
\begin{definition}
Such a formal sum is a \emph{chain} in $\Omega$ if each $\gamma_i^*\subset\Omega$.
It is a \emph{cycle} if it has a representation in which every
$\gamma_i$ is closed. Equivalently, the initial and terminal endpoints
in any representation can be paired. For the displayed representation,
write $\Gamma^*=\bigcup_i\gamma_i^*$.
\end{definition}
The sum is meaningful as addition of functionals, not as pointwise
addition of paths. The endpoint criterion follows from a combinatorial argument.
\src{199--200}
For a cycle and $\alpha\notin\Gamma^*$, set
\[
\Ind_\Gamma(\alpha)=\frac1{2\pi i}\int_\Gamma\frac{dz}{z-\alpha}
=\sum_i\Ind_{\gamma_i}(\alpha).
\]
Replacing each path by its opposite gives $-\Gamma$; then
$\int_{-\Gamma}f=-\int_\Gamma f$, and
$\Ind_{-\Gamma}(\alpha)=-\Ind_\Gamma(\alpha)$.
Chains can be added or subtracted through their functionals.
For integers $k_1,k_2$,
\[
\Gamma=k_1\Gamma_1\dotplus k_2\Gamma_2
\quad\Longrightarrow\quad
\int_\Gamma f=k_1\int_{\Gamma_1}f+k_2\int_{\Gamma_2}f
\]
for every $f$ continuous on $\Gamma_1^*\cup\Gamma_2^*$.
Chains can have many representations: equality of
$\sum_i\gamma_i$ and $\sum_j\delta_j$ means equality of
$\sum_i\int_{\gamma_i}f$ and $\sum_j\int_{\delta_j}f$
for every continuous $f$ on the union of all the path images.
In particular, a cycle can be represented by paths that are not closed.
\section{Global integral theorem and formula}
\src{201}
\begin{theorem}
Let $f\in H(\Omega)$ for an arbitrary open $\Omega\subseteq\CC$.
If a cycle $\Gamma$ in $\Omega$ satisfies
$\Ind_\Gamma(\alpha)=0$ for every $\alpha\notin\Omega$, then
\[
f(z)\Ind_\Gamma(z)=\frac1{2\pi i}\int_\Gamma\frac{f(w)}{w-z}\,dw
\quad(z\in\Omega\setminus\Gamma^*),\qquad \int_\Gamma f(z)\,dz=0.
\]
If cycles $\Gamma_0,\Gamma_1$ have equal indices at every
$\alpha\notin\Omega$, then $\int_{\Gamma_0}f=\int_{\Gamma_1}f$.
\end{theorem}
\src{206}
\begin{corollary}
For a cycle $\Gamma$ in an open set $\Omega$, the following are equivalent:
\begin{enumerate}
\item The preceding Cauchy integral formula holds for every $f\in H(\Omega)$.
\item $\int_\Gamma f(z)\,dz=0$ for every $f\in H(\Omega)$.
\item $\Ind_\Gamma(\alpha)=0$ for every $\alpha\in\CC\setminus\Omega$.
\end{enumerate}
\end{corollary}
\editorial{The universal quantifier on $f$ is essential. The slide
introduces one $f$ before the equivalence, but the assertion is false
for a single fixed function, for example $f=0$.}
\src{208--210}
\begin{remark}
In a convex region $\Omega$, apply the convex Cauchy theorem to
$(z-\alpha)^{-1}$, $\alpha\notin\Omega$, to obtain
$\Ind_\gamma(\alpha)=0$. Thus every cycle there satisfies the
global hypothesis, which generalizes the earlier convex results.
An efficient way to compute index is the winding number:
\[
W(\gamma,z_0)=\frac{\theta_{z_0}(b)-\theta_{z_0}(a)}{2\pi}
=\frac1{2\pi i}\int_\gamma\frac{dw}{w-z_0}
=\Ind_\gamma(z_0).
\]
It counts net turns. Each anticlockwise loop increases the argument
divided by $2\pi$ by one, and each clockwise loop decreases it by one.
The last part of the theorem specifies when one integration cycle
can be replaced by another without changing the integral.
\end{remark}
\src{211}
\begin{example}
For the illustrated path in $U$ and $f\in H(U)$,
$\Ind_\gamma(z_1)=\Ind_\gamma(z_2)=\Ind_\gamma(z_3)=1$.
Taking small anticlockwise circles about these points gives
$\int_\gamma f=\int_{\gamma_1}f+\int_{\gamma_2}f+\int_{\gamma_3}f$.
\sourcefigure{p211}{.72}{First cycle-deformation example.}
\end{example}
\src{212}
\begin{example}
In $\CC\setminus\{z_1,z_2,z_3\}$, the second illustrated path has
indices $1,2,2$. For positively oriented small circles and holomorphic $f$,
$\int_\gamma f=\int_{\gamma_1}f+2\int_{\gamma_2}f+2\int_{\gamma_3}f$.
\sourcefigure{p212}{.9}{Second cycle-deformation example.}
\end{example}
\src{213}
\begin{example}
In $\CC\setminus\{z_1,z_2,z_3,z_4\}$, the third path has indices
$-1,-2,-1,-2$. Using positively oriented small circles,
\[
\int_\gamma f=-\int_{\gamma_1}f-2\int_{\gamma_2}f
-\int_{\gamma_3}f-2\int_{\gamma_4}f.
\]
\sourcefigure{p213}{.95}{Third cycle-deformation example.}
\end{example}
\editorial{Slide 213 calls the small circles clockwise, but its arrows
and negative coefficients use anticlockwise circles. With clockwise
small circles, all four coefficients in the formula become positive.}
\src{214}
Recall the annular formula used for Laurent expansions.
\begin{theorem}
For $f$ analytic near $\overline A(z_0,r_1,r_2)$, and positively oriented
inner and outer circles $\gamma_1,\gamma_2$,
\[
f(z)=\frac1{2\pi i}\int_{\gamma_2}\frac{f(w)}{w-z}\,dw
-\frac1{2\pi i}\int_{\gamma_1}\frac{f(w)}{w-z}\,dw
\quad(r_1<|z-z_0|<r_2).
\]
\end{theorem}
\section{Homotopy and simple connectivity}
\src{215}
\begin{definition}
Closed curves $\gamma_0,\gamma_1:[0,1]\to X$ are $X$-\emph{homotopic}
if a continuous $H:[0,1]^2\to X$ satisfies
\[
H(s,0)=\gamma_0(s),\quad H(s,1)=\gamma_1(s),\quad H(0,t)=H(1,t).
\]
The curves $\gamma_t(s)=H(s,t)$ form a one-parameter family
continuously deforming one curve into the other inside $X$.
A curve homotopic to a constant curve is \emph{null-homotopic}.
\end{definition}
\src{216}
\begin{definition}
A connected space $X$ is \emph{simply connected} if every closed
curve in $X$ is null-homotopic.
\end{definition}
\begin{example}
Every convex region is simply connected: fix $z_1\in\Omega$
and set $H(s,t)=(1-t)\gamma_0(s)+tz_1$.
\end{example}
\begin{lemma}
For closed paths $\gamma_0,\gamma_1:[0,1]\to\CC$, if
$|\gamma_1(s)-\gamma_0(s)|<|\alpha-\gamma_0(s)|$ for all $s$,
then $\Ind_{\gamma_1}(\alpha)=\Ind_{\gamma_0}(\alpha)$.
\end{lemma}
\src{220}
\begin{theorem}
$\Omega$-homotopic closed paths in a region $\Omega$ have equal
indices at every point outside $\Omega$.
\end{theorem}
\src{224}
\begin{theorem}
In a simply connected region $\Omega$, every closed path $\Gamma$
has $\Ind_\Gamma(\alpha)=0$ for $\alpha\notin\Omega$.
For continuous closed curves, the same assertion uses winding number.
\end{theorem}
\src{225}
\begin{remark}
The global Cauchy theorem therefore applies in simply connected regions.
For $\Omega$-homotopic closed paths, equality of their outside indices
combined with the global theorem gives
$\int_{\Gamma_0}f(z)\,dz=\int_{\Gamma_1}f(z)\,dz$ for $f\in H(\Omega)$.
\end{remark}

\chapter{Analytic branches of the logarithm}
\lecturedetails{October 6, 2025}{226--248}
\section{Branches and arguments, recalled}
\src{227--228}
The exponential on $S_\alpha=\{x+iy:\alpha\le y<\alpha+2\pi\}$
is a bijection onto $\CC\setminus\{0\}$. Its inverse is $\log_\alpha$,
with $\arg_\alpha=\Im\log_\alpha$.
Thus $\log_\alpha(e^z)=z$ on this strip and $e^{\log_\alpha z}=z$ for $z\ne0$.
The principal branches are $\operatorname{Log}=\log_{-\pi}$ and
$\operatorname{Arg}=\arg_{-\pi}$.
\begin{theorem}
For $z\ne0$, $\log_\alpha z=\log|z|+i\arg_\alpha z$,
where $\arg_\alpha z\in[\alpha,\alpha+2\pi)$ is uniquely determined by
$z/|z|=e^{i\arg_\alpha z}$.
On the slit plane $\CC\setminus R_\alpha$, with
$R_\alpha=\{re^{i\alpha}:r\ge0\}$, both functions are continuous,
and $\log_\alpha$ is analytic with derivative $1/z$.
They are discontinuous at each nonzero point of the slit.
\end{theorem}
\src{230}
\begin{theorem}[Recalled inverse criterion]
If $g$ is analytic on $\Omega_1$, $g'\ne0$, $f$ is continuous on
open $\Omega$, $f(\Omega)\subseteq\Omega_1$, and $g(f(z))=z$,
then $f$ is analytic, one-to-one, and $f'=1/(g'\circ f)$.
\end{theorem}
\section{Analytic logarithms and powers}
\src{231}
\begin{definition}
For analytic $f$ on $\Omega$, an \emph{analytic logarithm} is an
analytic function $g$ on $\Omega$ such that $e^{g(z)}=f(z)$.
\end{definition}
\begin{remark}
The exponential is not injective on $\CC$, so logarithms require a
branch. Also, $f$ must be nonvanishing since an exponential is never zero.
On $A=\{z:1/2<|z|<2\}$ there is no analytic $g$ with $e^{g(z)}=z$:
differentiation would give $g'(z)e^{g(z)}=1$, hence $g'=1/z$,
but $1/z$ has no primitive on this annulus.
\end{remark}
\editorial{The slide concludes that the domain cannot have holes.
More precisely, holes can obstruct a logarithm of a particular function;
some functions, such as $f=1$, have logarithms on every domain.}
\src{232}
\begin{theorem}
If $\Omega$ is a simply connected region with $1\in\Omega$ and
$0\notin\Omega$, there is a branch $F=\log_\Omega$ which is holomorphic,
satisfies $e^{F(z)}=z$, and agrees with $\log r$ for positive real
$r$ sufficiently near $1$. Thus this branch extends the standard
real logarithm locally near $1$.
\end{theorem}
\src{235--236}
\sourcefigure{p235}{.78}{Principal branch of the logarithm on the slit plane (source slide 235).}
\begin{example}
On $\Omega=\CC\setminus(-\infty,0]$, write $z=re^{i\theta}$ with
$|\theta|<\pi$. The principal branch is $\log z=\log r+i\theta$.
The integration path consists of the real segment from $1$ to $r$
followed by the circular arc $\eta$ from $r$ to $z$:
\[
\log z=\int_1^r\frac{dx}{x}+\int_\eta\frac{dw}{w}
=\log r+\int_0^\theta\frac{ire^{it}}{re^{it}}\,dt
=\log r+i\theta.
\]
In general $\log(z_1z_2)\ne\log z_1+\log z_2$.
For $z_1=z_2=e^{2\pi i/3}$, their logarithms are $2\pi i/3$,
whereas $z_1z_2=e^{-2\pi i/3}$ and $\log(z_1z_2)=-2\pi i/3$.
\end{example}
\src{237}
For the principal branch, the Taylor expansion is
\[
\log(1+z)=z-\frac{z^2}{2}+\frac{z^3}{3}-\cdots
=\sum_{n=1}^\infty\frac{(-1)^{n+1}}n z^n,\qquad |z|<1.
\]
Both sides have derivative $1/(1+z)$ and value zero at zero.
\src{238}
On such a simply connected domain, choose $\log1=0$ and define,
for any $\alpha\in\CC$, $z^\alpha=e^{\alpha\log z}$.
Then $1^\alpha=1$. For $\alpha=1/n$,
\[
(z^{1/n})^n=\prod_{k=1}^n e^{(\log z)/n}
=e^{\sum_{k=1}^n(\log z)/n}=e^{\log z}=z.
\]
\src{239}
\begin{theorem}
Every nowhere-vanishing holomorphic function $f$ on a simply connected
region $\Omega$ has a holomorphic logarithm $g$, so that $f=e^g$.
The function $g$ may be denoted $\log f$ and determines a branch.
\end{theorem}
\src{241}
\begin{theorem}[Equivalent global Cauchy properties]
For a region $\Omega$, the following are equivalent:
\begin{enumerate}[label=(\Alph*)]
\item Every closed path $\gamma$ has $\Ind_\gamma(z)=0$ for $z\notin\Omega$.
\item $\int_\gamma f(w)\,dw=0$ for every closed path and every $f\in H(\Omega)$.
\item Every $f\in H(\Omega)$ has a primitive $F\in H(\Omega)$.
\item Every $f\in H(\Omega)$ with $1/f\in H(\Omega)$ can be written $f=e^g$
for some $g\in H(\Omega)$.
\item Every such $f$ can be written $f=g^2$ for some $g\in H(\Omega)$.
\end{enumerate}
\end{theorem}

% NEXT CHAPTER
\chapter{Meromorphic functions and residues}
\lecturedetails{October 9, 2025}{249--275}
\section{Meromorphic functions}
\src{250}
\begin{definition}
A function $f$ is meromorphic on an open set $\Omega$ if there is a set
$A\subset\Omega$ with no limit point in $\Omega$, such that
$f\in H(\Omega\setminus A)$ and $f$ has a pole at each point of $A$.
\end{definition}
\begin{remark}
$A=\varnothing$ is allowed, so every holomorphic function is meromorphic.
No compact subset contains infinitely many points of $A$; hence $A$ is
at most countable, since open subsets of $\CC$ are $\sigma$-compact.
\end{remark}
\src{251}
\begin{definition}
If $a$ is a pole and $Q(z)=\sum_{k=1}^m c_k(z-a)^{-k}$ is the principal
part (so $f-Q$ is removable at $a$), then
$\Res(f;a)=c_1$. For a cycle avoiding $a$,
\[
\frac1{2\pi i}\int_\Gamma Q(z)\,dz
=c_1\Ind_\Gamma(a).
\]
\end{definition}
\src{252}
\begin{theorem}[Residue theorem]
If $f$ is meromorphic on $\Omega$, $A$ is its pole set, and
$\Gamma$ is a cycle in $\Omega\setminus A$ with
$\Ind_\Gamma(\alpha)=0$ for $\alpha\notin\Omega$, then
\[
\frac1{2\pi i}\int_\Gamma f(z)\,dz
=\sum_{a\in A}\Res(f;a)\Ind_\Gamma(a).
\]
Only finitely many terms are nonzero.
\end{theorem}
\section{Argument principle and Rouche}
\src{254--257}
\begin{theorem}[Argument principle and Rouche's theorem]
Let $\gamma$ be a closed path in a region $\Omega$, with zero index
outside $\Omega$, and with index $0$ or $1$ on $\Omega\setminus\gamma^*$.
Let $\Omega_1=\{\alpha:\Ind_\gamma(\alpha)=1\}$, and let $N_f$ count
zeros in $\Omega_1$ with multiplicity.
If $f\in H(\Omega)$ has no boundary zero, then
\[
N_f=\frac1{2\pi i}\int_\gamma\frac{f'(z)}{f(z)}\,dz
=\Ind_{f\circ\gamma}(0).
\]
If $g\in H(\Omega)$ and $|f-g|<|f|$ on $\gamma^*$, then $N_g=N_f$.
\end{theorem}
\begin{remark}
Rouche's theorem preserves the number of zeros when two holomorphic
functions are sufficiently close on the boundary.
\end{remark}
\begin{example}
For $g(z)=z^{87}+36z^{57}+71z^4+z^3-z+1$, take $f(z)=71z^4$.
On $|z|=1$,
\[
|f-g|\le1+36+1+1+1<71=|f|.
\]
Thus $g$ has four zeros in $|z|<1$, counted with multiplicity.
\end{example}
\section{Contour examples}
\src{259--262}
\begin{example}
\[
\lim_{A\to\infty}\int_{-A}^{A}\frac{\sin x}{x}e^{itx}\,dx
=\begin{cases}\pi,&|t|<1,\\0,&|t|>1,\end{cases}
\]
and the value is $\pi/2$ when $t=\pm1$. Using
$2i\sin z=e^{iz}-e^{-iz}$ on the indented contour $\Gamma_A$ gives
\[
\int_{\Gamma_A}\frac{\sin z}{z}e^{itz}\,dz
=\phi_A(t+1)-\phi_A(t-1),
\quad
\phi_A(s)=\frac1{2\pi i}\int_{\Gamma_A}\frac{e^{isz}}z\,dz.
\]
Closing in the lower and upper half-planes gives
\[
\phi_A(s)=\frac1{2\pi}\int_{-\pi}^0e^{isAe^{i\theta}}d\theta,
\qquad
\phi_A(s)=1-\frac1{2\pi}\int_0^\pi e^{isAe^{i\theta}}d\theta.
\]
Since $e^{isAe^{i\theta}}=e^{-As\sin\theta}$, dominated convergence yields
$\phi_A(s)\to0$ for $s<0$ and $\to\pi$ for $s>0$; also
$\phi_A(0)=\pi/2$.
\end{example}
\src{263--265}
\begin{example}
\[
\int_{-\infty}^{\infty}\frac{dx}{1+x^2}=\pi.
\]
For $f(z)=(1+z^2)^{-1}$, the upper semicircle encloses $i$ and
\[
f(z)=\frac1{2i(z-i)}-\frac1{2i(z+i)},\qquad \Res(f;i)=\frac1{2i}.
\]
The contour integral is $\pi$, while the arc is bounded by $M/R\to0$.
\end{example}
\src{266--269}
\begin{example}
For $0<a<1$,
\[
\int_{-\infty}^{\infty}\frac{e^{ax}}{1+e^x}\,dx=\frac{\pi}{\sin\pi a}.
\]
Use $f(z)=e^{az}/(1+e^z)$ on a rectangle of height $2\pi$.
The sole pole is $\pi i$, with residue $-e^{a\pi i}$. If $I_R$ is
the bottom integral and the vertical sides vanish, the top is
$-e^{2\pi ia}I_R$, so
\[
(1-e^{2\pi ia})I=-2\pi i e^{a\pi i}
\quad\Longrightarrow\quad I=\frac{\pi}{\sin\pi a}.
\]
\end{example}
\src{270--275}
\begin{example}
\[
\int_{-\infty}^{\infty}\frac{e^{-2\pi ix\xi}}{\cosh\pi x}\,dx
=\frac1{\cosh\pi\xi}.
\]
Here $\cosh z=(e^z+e^{-z})/2$.
Use $f(z)=e^{-2\pi iz\xi}/\cosh\pi z$ on a rectangle of height $2i$.
The poles are $\alpha=i/2$ and $\beta=3i/2$, with residues
$e^{\pi\xi}/(\pi i)$ and $-e^{3\pi\xi}/(\pi i)$. The vertical sides
vanish, the top is $-e^{4\pi\xi}I$, and the residue theorem gives
$I=1/\cosh(\pi\xi)$.
\end{example}
\begin{remark}
The same method gives
\[
\int_\RR\frac{e^{-2\pi ix\xi}}{\cosh\pi x+\cos\pi a}\,dx
=\frac{\sin\pi a}{\sinh(2\pi\xi)}
\]
in the source's stated parameter range, with
$\sinh z=(e^z-e^{-z})/2$.
\end{remark}

\chapter{Phragmen--Lindelof and harmonic analysis}
\lecturedetails{October 13, 2025}{276--303}
\section{Phragmen--Lindelof method}
\src{277--280}
\begin{remark}
The maximum modulus principle fails on unbounded sets. On
$\Omega=\{x+iy:-\pi/2<y<\pi/2\}$, $f(z)=\exp(\exp z)$ has boundary
modulus $1$ but $f(x)=e^{e^x}\to\infty$. Growth control replaces
boundedness in Phragmen--Lindelof arguments. For example, an entire
function satisfying $|f(z)|\le1+|z|^{1/2}$ is constant.
\end{remark}
\begin{theorem}[Strip Phragmen--Lindelof]
Let $\Omega=\{x+iy:a<x<b\}$, with $f$ continuous and holomorphic
and $|f|<B$. Set $M(x)=\sup_y|f(x+iy)|$. Then
\[
M(x)^{b-a}\le M(a)^{b-x}M(b)^{x-a},\qquad a\le x\le b.
\]
\end{theorem}
\src{287}
\begin{corollary}
If the preceding assumptions hold and $f$ is nonconstant, then
$|f(z)|<\max(M(a),M(b))$ in $\Omega$.
\end{corollary}
\src{288}
\begin{theorem}[Hadamard three-circle theorem]
For $0<R_1<R_2$, let $g$ be holomorphic on the annulus and
$B(R)=\max_{|z|=R}|g(z)|$. Then
\[
B(R)^{\log R_2-\log R_1}
\le B(R_1)^{\log R_2-\log R}B(R_2)^{\log R-\log R_1}.
\]
\end{theorem}
\src{291}
\begin{theorem}[Strip application]
If $|\Im z|<\pi/2$, $f$ is holomorphic and continuous there, and
$|f(x+iy)|<\exp(Ae^{\alpha|x|})$ for $\alpha<1$, while
$|f(x\pm\pi i/2)|\le1$, then $|f(z)|\le1$ throughout the strip.
\end{theorem}
\begin{remark}
The conclusion fails at the critical growth $\alpha=1$, as
$\exp(\exp z)$ demonstrates.
\end{remark}
\section{Interpolation and convolution}
\src{294--301}
\begin{theorem}[Riesz interpolation]
For $\sigma$-finite measure spaces $(X,\mu),(Y,\nu)$, let $T$ be linear
on finitely simple functions and suppose
\[
\|Tf\|_{q_j}\le M_j\|f\|_{p_j},\quad j=0,1.
\]
For $0<\theta<1$, define
\[
\frac1p=\frac{1-\theta}{p_0}+\frac{\theta}{p_1},\qquad
\frac1q=\frac{1-\theta}{q_0}+\frac{\theta}{q_1}.
\]
Then $\|Tf\|_q\le M_0^{1-\theta}M_1^\theta\|f\|_p$.
\end{theorem}
\begin{remark}
Finitely simple functions are dense in $L^p$ for $p<\infty$ on
$\sigma$-finite spaces, yielding a unique bounded extension.
\end{remark}
\begin{lemma}[Three-lines lemma]
If $F$ is analytic in $0<\Re z<1$, bounded and continuous on the
closure, with $|F|\le B_0$ on $\Re z=0$ and $|F|\le B_1$ on $\Re z=1$,
then $|F(z)|\le B_0^{1-\theta}B_1^\theta$ on $\Re z=\theta$.
\end{lemma}
\begin{exercise}
Prove the three-lines lemma.
\end{exercise}
\src{302--303}
\begin{theorem}[Young's convolution inequality]
If $1\le p,q,r\le\infty$ and $1+1/q=1/p+1/r$, then
\[
\|f*g\|_{L^q(\RR)}\le\|f\|_{L^p(\RR)}\|g\|_{L^r(\RR)},
\quad
(f*g)(x)=\int_\RR f(x-y)g(y)\,dy.
\]
\end{theorem}
\begin{exercise}
For fixed $g\in L^r$, verify directly that convolution maps
$L^1$ to $L^r$ and $L^{r'}$ to $L^\infty$ with the stated norms,
then derive Young's inequality by interpolation.
\end{exercise}

\chapter{The Fourier transform}
\lecturedetails{October 16, 2025}{304--327}
\src{305--307}
\begin{definition}
For $a>0$, let $\mathcal F_a$ be the class of functions holomorphic on
$S_a=\{z:|\Im z|<a\}$ and satisfying a uniform exponential bound
of the form $|f(x+iy)|\le C e^{-b|x|}$ on every narrower strip
$|\Im z|\le b<a$. Such bounds imply the corresponding bounds for all
derivatives by Cauchy's formula.
\end{definition}
\src{308--311}
\begin{theorem}
If $f\in\mathcal F_a$, then for each $0<b<a$ there is $B$ such that
\[
|\widehat f(\xi)|\le B e^{-2\pi b|\xi|},\qquad
\widehat f(\xi)=\int_\RR f(x)e^{-2\pi ix\xi}\,dx.
\]
\end{theorem}
The proof shifts the real contour to $\Im z=-b$ for $\xi>0$
and to $\Im z=b$ for $\xi<0$; the vertical sides vanish by the
exponential estimates.
\src{312--318}
\begin{remark}
The transform has rapid decay. The Fourier inversion formula is
\[
f(x)=\int_\RR\widehat f(\xi)e^{2\pi ix\xi}\,d\xi
\]
under the source's decay hypotheses. The auxiliary identity used is
\[
\int_\RR\frac{\sin(A\xi)}{\xi}\,d\xi
=\pi\,\operatorname{sgn}(A).
\]
\end{remark}
\begin{theorem}[Poisson summation]
For $f$ in the source class,
\[
\sum_{n\in\ZZ}f(n)=\sum_{n\in\ZZ}\widehat f(n).
\]
\end{theorem}
\src{323--327}
\begin{remark}
The Gaussian identity gives the theta transformation. For $t>0$,
\[
\vartheta(t)=\sum_{n\in\ZZ}e^{-\pi tn^2},\qquad
\vartheta(t)=t^{-1/2}\vartheta(1/t).
\]
\end{remark}
\begin{theorem}[Paley--Wiener type statement]
If $|\widehat f(\xi)|\le Ae^{-2\pi a|\xi|}$, then
\[
f(z)=\int_\RR\widehat f(\xi)e^{2\pi i\xi z}\,d\xi
\]
defines a holomorphic function in the strip $|\Im z|<a$.
\end{theorem}

\chapter{Conformal mappings}
\lecturedetails{October 20, 2025}{328--359}
\src{329--331}
\begin{definition}
A conformal map $\Omega_1\to\Omega_2$ is a one-to-one holomorphic map
onto $\Omega_2$ with nonzero derivative. It is an analytic homeomorphism,
preserving oriented angles and infinitesimal shapes.
\end{definition}
\src{332--339}
\begin{theorem}
The Cayley map
\[
F(z)=\frac{i-z}{i+z}
\]
conformally maps the upper half-plane $\HH$ onto the unit disc $\DD$.
\end{theorem}
\begin{theorem}
For $n\ge1$, the map $z\mapsto z^n$ conformally maps
$\{0<\arg z<\pi/n\}$ onto $\HH$.
\end{theorem}
\begin{theorem}
The map $z\mapsto z^2$ conformally maps the upper half-disc onto
the first quadrant. The exponential maps $\HH$ conformally onto
the horizontal strip $\{0<\Im z<\pi\}$.
\end{theorem}
\src{340--350}
\begin{theorem}[Schwarz lemma]
If $f\in H(\DD)$, $|f|\le1$, and $f(0)=0$, then $|f(z)|\le|z|$.
If equality holds at a nonzero point, $f(z)=\lambda z$, $|\lambda|=1$.
\end{theorem}
\begin{definition}
\[
\varphi_a(z)=\frac{a-z}{1-\bar az},\qquad a\in\DD.
\]
\end{definition}
\begin{remark}
The maps $\varphi_a$ are disc automorphisms, satisfy
$\varphi_a^{-1}=\varphi_a$, and preserve the disc. Schwarz's lemma
applied after these automorphisms gives
\[
|f'(w)|\le\frac{1-|f(w)|^2}{1-|w|^2}.
\]
\end{remark}
\begin{theorem}
If $f$ is an automorphism of $\DD$ and $f(w)=0$, then
$f=\lambda\varphi_w$ for some $|\lambda|=1$.
\end{theorem}
\src{351--359}
\begin{definition}
\[
\mathrm{SL}_2(\RR)=\left\{\begin{pmatrix}a&b\\c&d\end{pmatrix}:ad-bc=1\right\},
\qquad f_M(z)=\frac{az+b}{cz+d}.
\]
\end{definition}
\begin{theorem}
$\Aut(\HH)=\{f_M:M\in\mathrm{SL}_2(\RR)\}$. The Cayley map conjugates
these transformations to the automorphisms of $\DD$.
\end{theorem}

\chapter{The Riemann mapping theorem}
\lecturedetails{October 23, 2025}{360--386}
\begin{theorem}[Riemann mapping]
Every simply connected region $\Omega\ne\CC$ is conformally equivalent
to $\DD$.
\end{theorem}
\src{362--375}
\begin{lemma}
Every open set $\Omega\subset\CC$ admits compact sets $K_n\subset\Omega$
with $K_n\subset\operatorname{int}K_{n+1}$ and $\bigcup_nK_n=\Omega$.
\end{lemma}
\begin{lemma}[Hurwitz]
If $f_n\in H(\Omega)$ converge uniformly on compact subsets to $f$,
and each $f_n$ is zero-free, then $f$ is zero-free or identically zero.
\end{lemma}
\begin{definition}
A family $\mathcal F\subset H(\Omega)$ is normal if every sequence in
it has a subsequence converging uniformly on compact subsets to a
holomorphic function or to infinity.
\end{definition}
\begin{theorem}[Montel]
Every locally uniformly bounded family of holomorphic functions is normal.
\end{theorem}
\src{376--386}
The proof of the mapping theorem uses the family
\[
\Sigma=\{\psi\in H(\Omega):\psi\text{ one-to-one and }\psi(\Omega)\subset\DD\},
\]
normality, maximization of $|\psi'(z_0)|$, disc automorphisms,
and Schwarz's lemma. The extremal map is onto $\DD$.

\chapter{Runge's theorem and approximation}
\lecturedetails{October 27, 2025}{387--416}
\src{389--397}
\begin{definition}
A rational function is $P/Q$ for polynomials $P,Q$, $Q\not\equiv0$.
Its poles are the zeros of $Q$ after cancellation.
\end{definition}
\begin{theorem}[Cauchy theorem for compact sets]
For a compact $K$ and an open neighborhood $\Omega$, a finite balanced
collection of oriented grid intervals gives a cycle whose index is one
on $K$ and zero outside $\Omega$.
\end{theorem}
\begin{theorem}[Runge]
Let $K\subset\CC$ be compact and choose one point $\alpha_j$ in each
component of $\widehat\CC\setminus K$. If $\Omega\supset K$ is open and
$f\in H(\Omega)$, then $f$ is uniformly approximated on $K$ by rational
functions with poles among the $\alpha_j$.
\end{theorem}
\src{403--408}
\begin{corollary}[Polynomial approximation]
If $\widehat\CC\setminus K$ is connected, every function holomorphic
near $K$ is uniformly approximated on $K$ by polynomials.
\end{corollary}
\begin{remark}
Connectedness of the complement is necessary: holes require poles
in the complementary components.
\end{remark}
\src{409--411}
\begin{theorem}[Mittag--Leffler]
Given a discrete set of points and prescribed principal parts there,
there exists a meromorphic function with exactly those principal parts.
\end{theorem}
\src{412--416}
\begin{theorem}[Characterizations of simple connectivity]
For a region $\Omega$, the following conditions are equivalent:
\begin{enumerate}
\item $\Omega$ is conformally equivalent to $\DD$.
\item $\Omega$ is simply connected.
\item Every holomorphic function on $\Omega$ has a primitive.
\item Every nonvanishing holomorphic function has a logarithm.
\item Every nonvanishing holomorphic function has a holomorphic square root.
\item Every holomorphic function near compact subsets is approximable by polynomials.
\end{enumerate}
\end{theorem}

\chapter{Harmonic functions}
\lecturedetails{October 30, 2025}{417--456}
\begin{theorem}[Cauchy--Riemann in Wirtinger form]
\[
\partial_x f=\frac1i\,\partial_y f.
\]
\end{theorem}
\begin{definition}
A real function $u$ is harmonic when $u_{xx}+u_{yy}=0$.
\end{definition}
\begin{theorem}
The real and imaginary parts of a holomorphic function are harmonic.
A harmonic function vanishing on a nonempty open subset of a region
vanishes throughout that region.
\end{theorem}
\begin{definition}
If $u$ is harmonic, $v$ is a harmonic conjugate when $u+iv$ is holomorphic;
equivalently $u_x=v_y$ and $u_y=-v_x$.
\end{definition}
\begin{theorem}
Every harmonic function on a disc has a harmonic conjugate. A region is
simply connected exactly when every harmonic function on it has one.
\end{theorem}
\begin{definition}
Continuous $u$ has the mean-value property if
\[
u(a)=\frac1{2\pi}\int_0^{2\pi}u(a+re^{i\theta})\,d\theta
\]
whenever the closed disc is contained in the region.
\end{definition}
\begin{theorem}
An $MVP$ function satisfies the maximum principle; a nonconstant one
cannot attain an interior maximum or minimum. Conversely, a continuous
$MVP$ function is harmonic.
\end{theorem}
\begin{definition}
\[
P_r(\theta)=\Re\frac{1+re^{i\theta}}{1-re^{i\theta}}
=\frac{1-r^2}{1-2r\cos\theta+r^2}.
\]
\end{definition}
\begin{theorem}[Dirichlet problem for a disc]
For continuous real boundary data $f$ on $C(a,\rho)$ there is a unique
continuous harmonic $u$ on the closed disc with $u=f$ on its boundary,
given by the Poisson integral
\[
u(a+re^{i\theta})=\frac1{2\pi}\int_{-\pi}^{\pi}
P_{r/\rho}(\theta-\varphi)f(a+\rho e^{i\varphi})\,d\varphi.
\]
\end{theorem}

\chapter{Infinite products}
\lecturedetails{November 3, 2025}{457--494}
\begin{theorem}[Jensen]
If $f$ is holomorphic near $\overline{D(0,R)}$, $f(0)\ne0$, and
$a_1,\ldots,a_N$ are its zeros in the disc, counted with multiplicity,
\[
\log|f(0)|+\sum_{j=1}^N\log\frac R{|a_j|}
=\frac1{2\pi}\int_0^{2\pi}\log|f(Re^{i\theta})|\,d\theta.
\]
\end{theorem}
\begin{definition}
An entire function has order at most $\rho$ if
$|f(z)|\le Ae^{B|z|^\rho}$ for suitable $A,B$.
\end{definition}
\begin{theorem}
An entire function of order at most $\rho$ has at most $Cr^\rho$
zeros in $|z|\le r$ for large $r$.
\end{theorem}
\begin{definition}
An infinite product $\prod(1+a_n)$ converges if its partial products
converge to a nonzero limit. If $\sum|a_n|<\infty$, it does.
\end{definition}
\begin{lemma}
If $|z|<1/2$, then $|\log(1+z)|\le2|z|$.
\end{lemma}
\begin{theorem}[Weierstrass product]
For every sequence $|a_n|\to\infty$, an entire function exists whose
zeros are exactly the prescribed $a_n$, with prescribed multiplicities.
\end{theorem}
\begin{definition}
\[
E_0(z)=1-z,\qquad
E_k(z)=(1-z)\exp\left(z+\frac{z^2}{2}+\cdots+\frac{z^k}{k}\right).
\]
\end{definition}
\begin{theorem}[Hadamard factorization]
An entire function of finite order factors as a polynomial times
canonical factors $E_k(z/a_n)$ and an exponential factor of bounded degree.
\end{theorem}

\chapter{Gamma function and summation}
\lecturedetails{November 10, 2025}{495--549}
\begin{theorem}
\[
\sin\pi z=\pi z\prod_{n=1}^{\infty}\left(1-\frac{z^2}{n^2}\right).
\]
\end{theorem}
\begin{definition}
\[
\Gamma(z)=\frac{e^{-\gamma z}}{z}\prod_{n=1}^{\infty}
\left(1+\frac zn\right)^{-1}e^{z/n}.
\]
\end{definition}
\begin{theorem}[Euler's functional equation]
\[
\Gamma(z+1)=z\Gamma(z),\qquad\Gamma(n+1)=n!.
\]
\end{theorem}
\begin{theorem}[Reflection and duplication]
\[
\Gamma(z)\Gamma(1-z)=\frac{\pi}{\sin\pi z},\qquad
\Gamma(2z)\Gamma(1/2)=2^{2z-1}\Gamma(z)\Gamma(z+1/2).
\]
\end{theorem}
\begin{theorem}[Integral representation]
For $\Re z>0$, $\Gamma(z)=\int_0^\infty e^{-t}t^{z-1}\,dt$.
\end{theorem}
\begin{definition}
\[
\frac{ze^{xz}}{e^z-1}=\sum_{k=0}^{\infty}B_k(x)\frac{z^k}{k!}.
\]
\end{definition}
\begin{theorem}[Euler--Maclaurin]
\[
\sum_{n=a}^{b}f(n)=\int_a^bf(x)\,dx+\frac{f(a)+f(b)}2+
\sum_{r=2}^{q}\frac{B_r}{r!}(f^{(r-1)}(b)-f^{(r-1)}(a))+R_q.
\]
\end{theorem}
\begin{theorem}[Stirling]
\[
\Gamma(z)=\sqrt{2\pi}\,z^{z-1/2}e^{-z}(1+O(1/z))
\]
in sectors $|\arg z|\le\pi-\delta$.
\end{theorem}

\chapter{Dirichlet series}
\lecturedetails{November 13, 2025}{550--594}
\begin{definition}
\[
D(s)=\sum_{n=1}^{\infty}a_n n^{-s},\qquad s=\sigma+it.
\]
\end{definition}
\begin{theorem}
There are abscissae $\sigma_a$ and $\sigma_c$ of absolute convergence
and convergence. In every strict half-plane $\sigma>\sigma_c$, $D$
is holomorphic and may be differentiated term by term.
\end{theorem}
\begin{lemma}
The coefficients of a Dirichlet series are uniquely determined by the
function in a half-plane of convergence.
\end{lemma}
\begin{definition}
A function $f$ is multiplicative when $f(mn)=f(m)f(n)$ for coprime $m,n$.
\end{definition}
\begin{theorem}[Euler product]
\[
\sum_{n\ge1}\frac{f(n)}{n^s}
=\prod_p\left(\sum_{k\ge0}\frac{f(p^k)}{p^{ks}}\right)
\]
under the source's convergence hypotheses.
\end{theorem}
\begin{theorem}
Products of Dirichlet series correspond to Dirichlet convolution of
their coefficients, with absolute convergence justifying rearrangement.
\end{theorem}

\chapter{The zeta function and the prime number theorem}
\lecturedetails{November 17, 2025}{595--631}
\begin{definition}
\[
\Lambda(n)=\begin{cases}\log p,&n=p^k,\\0,&\text{otherwise},\end{cases}
\quad\psi(x)=\sum_{n\le x}\Lambda(n),\quad
\vartheta(x)=\sum_{p\le x}\log p.
\]
\end{definition}
\begin{theorem}[Chebyshev bounds]
\[
\vartheta(x)\le x\log4,\qquad\psi(x)\asymp x.
\]
\end{theorem}
\begin{definition}
\[
\zeta(s)=\sum_{n=1}^{\infty}n^{-s}
=\prod_p(1-p^{-s})^{-1}\quad(\Re s>1).
\]
\end{definition}
\begin{theorem}[Analytic continuation and functional equation]
\[
\xi(s)=s(s-1)\pi^{-s/2}\Gamma(s/2)\zeta(s)
\]
extends to an entire function and satisfies $\xi(s)=\xi(1-s)$.
\end{theorem}
\begin{remark}
The trivial zeros are $-2,-4,-6,\ldots$; $\zeta$ has no zeros for
$\Re s>1$, and the nontrivial zeros lie in $0\le\Re s\le1$.
\end{remark}
\begin{theorem}[Explicit formula]
\[
\psi(x)=x-\sum_\rho\frac{x^\rho}{\rho}
-\log(2\pi)-\frac12\log(1-x^{-2})
\text{the source's limiting correction}.
\]
\end{theorem}
\begin{theorem}[Prime number theorem]
\[
\psi(x)=x+O\!\left(xe^{-c\sqrt{\log x}}\right).
\]
\end{theorem}
\begin{conjecture}[Riemann hypothesis]
Every nontrivial zero of $\zeta$ has real part $1/2$.
\end{conjecture}
\begin{theorem}[Riemann--von Mangoldt]
\[
N(T)=\frac{T}{2\pi}\log\frac{T}{2\pi}-\frac{T}{2\pi}+O(\log T).
\]
\end{theorem}

\chapter{Quantitative estimates for the prime number theorem}
\lecturedetails{November 20, 2025}{632--693}
\begin{lemma}
The logarithmic derivative $\zeta'/\zeta$ admits a renormalized
partial-fraction expansion in terms of the zeros, with absolute
convergence in bounded vertical strips away from zeros and gamma poles.
\end{lemma}
\begin{theorem}[Zero-free region]
There is $C>0$ such that $\zeta(s)$ has no zero
\[
\rho=\beta+i\gamma,\qquad
\beta\ge1-\frac{C}{\log(|\gamma|+3)}.
\]
\end{theorem}
\begin{lemma}[Perron kernel]
\[
\frac1{2\pi i}\int_{\kappa-iT}^{\kappa+iT}\frac{x^s}{s}\,ds
\]
is $1$ for $x>1$, $0$ for $0<x<1$, and has the half-value at $x=1$,
up to an error controlled by $x^\kappa/(1+T|\log x|)$.
\end{lemma}
\begin{theorem}[Truncated Perron formula]
For a Dirichlet series, the finite-height Perron integral recovers
the summatory function with the explicit truncation error displayed
in the source lecture.
\end{theorem}
\begin{theorem}[Landau explicit formula]
\[
\psi(x)=x-\sum_{|\gamma|\le T}\frac{x^\rho}{\rho}
-\log(2\pi)-\frac12\log(1-x^{-2})+\text{truncation error}.
\]
\end{theorem}
\begin{theorem}[Korobov--Vinogradov]
There is $C>0$ such that no zero lies in
\[
\sigma\ge1-\frac{C}{(\log(|t|+3))^{2/3}
(\log\log(|t|+3))^{1/3}}.
\]
\end{theorem}
\begin{remark}
The source obtains its best stated PNT error term by inserting this
zero-free region into Perron's formula.
\end{remark}

\end{document}
